% ============================================================
%  第二卷 第十二章 引力物体的场（§99–§106 + 习题）
%  底本：高教社中文版第8版；OCR 错误已修正；脚注文本已由脚注代理补录（2026-10-10）
%  主要修正：
%  - §99 为本章首节（原书目录 §99 与章名同起于 p.329），非第十一章末尾；
%  - OCR 将克里斯托费尔符号 Γ 排为 T/𝒯，统一恢复为 \Gamma；
%  - OCR ℓ0/δ′ 为 𝔈0/𝔈′ 之误（§101 两处、§102 习题1 一处、§99 无）；
%  - §99：Γ_{00}^α 式与 R_0^0 定义各补回一负号（与 (99.2) 自洽）；
%    "我们省略了μ的下标0" OCR 误为上标 0)；
%  - §100：(100.16) 前一行 e^{-λ}/τ → e^{-λ}/r；习题1 R_{2323} = -r_g sin²θ
%    → -r r_g sin²θ（用不变量 I₁=(r_g/2r³)² 核对）；
%  - §102：首句"史瓦西度规(101.4)"应为 (100.14)；习题1 r_g=2km′/c 漏平方；
%    "收敛物体的半径 r_g" 衍文 r_g 删；句末页码残迹 66 删；
%  - §103：(103.8)—(103.10) 条件"当 f>0/f<0/f=0"乱码复原；m 积分下限 r→0；
%    "导数是其径向速度"补 ṙ；(103.15) 尾部 OCR 乱码剔除；
%  - §104：(104.1) 2kM/r → 2kM/(c²r)（与 §105 习题1 式(1)一致）；
%    (106.9) 式后正文 m_α → m_a；
%  - §106：(106.17) 烂尾分式按原书补全 k²m_am_bm_c/(2m_a r_ab r_ac)；
%    "形如(105.6a)"→(105.6)；习题4 δ^{(1)}L 积分系数 3km′/2c² → 3km′/c²
%    （|V+ω×r|² 展开因子 2，与末式 Ω 自洽）；
%  - 习题内公式保留原书局部编号 (1)(2)…（以 \tag 实现）；正文公式自动编号；
%  - 脚注已全部转录（48 处圈码 + §104 星号脚注 1 处 → \footnote{}），逐条对照
%    原书扫描页页脚录入（原书转换/pages/p-NNN.png，PDF 页 = 印张页 + 14）；
%  - §101“光线偏折”一句的脚注（原书 p.345 底，δφ=1.75″）OCR 漏识别圈码，
%    已按原书补加 \footnote；脚注由 footmisc(perpage) 按页自动编号为圈码；
%  - 插图：图20（§102）、图21/22/23（§102 习题）、图24（§103）。
% ============================================================
\chapter{引力物体的场}

\section{牛顿定律}
% 注：§99 按原书目录属第十二章首节（p.329），非第十一章末尾，未外移。

现在，我们将爱因斯坦场方程过渡到非相对论力学极限.正如在§87所指出的，关于所有粒子速度很小的假设也要求引力场很弱.

在§87中，我们求得在所考虑的极限情形下，度规张量的分量 $g_{00}$（我们所需要的唯一分量）是
\begin{equation*}
  g_{00}=1+\frac{2\varphi}{c^{2}}.
\end{equation*}

其次，对于能量动量张量的分量，我们可以用表达式 (35.4) $T_i^k=\mu c^2u_iu^k$，其中 $\mu$ 是物体的密度（单位体积内粒子静止质量之和；我们省略了 $\mu$ 的下标 0）.因为宏观运动也认为是缓慢的，那么，对于四维速度 $u^i$，我们应当略去它的所有的空间分量，而仅仅保留时间分量，就是说，应当令 $u^{\alpha}=0,\ u^0=u_0=1$.因此，在 $T_i^k$ 的所有分量之中只留有
\begin{equation}
  T_0^0=\mu c^{2}.
\end{equation}

标量 $T=T_i^i$ 将等于同一值 $\mu c^{2}$.

我们将爱因斯坦方程 (95.8) 写成
\begin{equation*}
  R_i^k=\frac{8\pi k}{c^{4}}\left(T_i^k-\frac{1}{2}\delta_i^kT\right);
\end{equation*}
当 $i=k=0$ 时，
\begin{equation*}
  R_0^0=\frac{4\pi k}{c^{2}}\mu.
\end{equation*}

不难验证，在我们所考虑的近似情形中，所有其余的方程都恒等于零.

在从普遍公式 (92.7) 计算 $R_0^0$ 时，我们要注意，所有含 $\Gamma_{kl}^i$ 诸量的乘积的项，在任何情形下，都是二阶小量，含有对于 $x^0=ct$ 的导数的项是很小的（同那些含有对于坐标 $x^{\alpha}$ 的导数的项比较），因为它们含有额外的 $1/c$ 的幂.结果，剩下 $R_0^0=R_{00}=-\partial\Gamma_{00}^{\alpha}/\partial x^{\alpha}$.将
\begin{equation*}
  \Gamma_{00}^{\alpha}\approx-\frac{1}{2}g^{\alpha\beta}
  \frac{\partial g_{00}}{\partial x^{\beta}}
  =-\frac{1}{c^{2}}\frac{\partial\varphi}{\partial x^{\alpha}}
\end{equation*}
代入，我们求得
\begin{equation*}
  R_0^0=\frac{1}{c^{2}}\frac{\partial^{2}\varphi}{\partial x^{\alpha 2}}
  \equiv\frac{1}{c^{2}}\Delta\varphi.
\end{equation*}

因此，爱因斯坦方程化为
\begin{equation}
  \Delta\varphi=4\pi k\mu.
\end{equation}

这就是非相对论力学中的引力场方程.我们要注意这个方程完全类似于电势的泊松方程 (36.4)，只是在此处以质量密度乘以 $k$ 来代替了电荷密度.因此，与 (36.8) 相类似，我们可以立即写出方程 (99.2) 的通解如下：
\begin{equation}
  \varphi=-k\int\frac{\mu}{R}\,\mathrm{d}V.
\end{equation}

这个公式决定了非相对论近似下任意质量分布的引力场的势.

就特例言之，质量为 $m$ 的单个粒子的场的势是
\begin{equation}
  \varphi=-\frac{km}{R},
\end{equation}

因此，作用在该场内另外一个粒子（质量为 $m'$）上的力 $F=-m'\partial\varphi/\partial R$ 等于
\begin{equation}
  F=-\frac{kmm'}{R^{2}}.
\end{equation}

这就是众所周知的牛顿引力定律.

粒子在引力场内的势能等于粒子的质量乘以场的势，这与电场内的势能等于电荷乘以场的势相似.因此，与 (37.1) 相似，我们可以写出任意质量分布的势能如下：
\begin{equation}
  U=\frac{1}{2}\int\mu\varphi\,\mathrm{d}V.
\end{equation}

对于离产生场的质量很远的恒定引力场的牛顿势，我们可以写出一个类似于在§40和§41中对静电场得到的展开式.选择质量的惯性中心为坐标原点.这时，积分 $\int\mu\boldsymbol{r}\,\mathrm{d}V$ 将恒等于零，这个积分与电荷体系的偶极矩类似.因此，与静电场的情形不同，在引力场情形中，我们总能够消除“偶极项”.因此，势 $\varphi$ 的展开式有下面的形式：
\begin{equation}
  \varphi=-k\left(\frac{M}{R_0}
  +\frac{1}{6}D_{\alpha\beta}
  \frac{\partial^{2}}{\partial X_{\alpha}\partial X_{\beta}}
  \frac{1}{R_0}+\cdots\right),
\end{equation}
式中 $M=\int\mu\,\mathrm{d}V$ 是体系的总质量，而
\begin{equation}
  D_{\alpha\beta}=\int\mu(3x_{\alpha}x_{\beta}-r^{2}\delta_{\alpha\beta})
  \,\mathrm{d}V
\end{equation}
可以称为质量的四极矩张量\footnote{在此，我们将所有指标 $\alpha,\beta$ 都写成下标，没有区分协变和逆变分量，因为所有的运算都是在普通的牛顿空间（欧几里得空间）中进行的.}.它与通常的转动惯量张量
\begin{equation*}
  J_{\alpha\beta}=\int\mu(r^{2}\delta_{\alpha\beta}-x_{\alpha}x_{\beta})
  \,\mathrm{d}V
\end{equation*}
有下面的关系式：
\begin{equation}
  D_{\alpha\beta}=J_{\gamma\gamma}\delta_{\alpha\beta}-3J_{\alpha\beta}.
\end{equation}

由给定质量分布求牛顿势是数学物理中一个分支的主题；阐述有关的各种方法并不是本书的任务.这里我们只给出一个均匀椭球体产生的引力场势的公式作为参考.

设椭球表面由如下方程给出：
\begin{equation}
  \frac{x^{2}}{a^{2}}+\frac{y^{2}}{b^{2}}+\frac{z^{2}}{c^{2}}=1,
  \qquad a>b>c.
\end{equation}

则场在物体外任一点的势由如下公式给定：
\begin{equation}
\begin{aligned}
  &\varphi=-\pi\mu abck\int_{\xi}^{\infty}
  \left(1-\frac{x^{2}}{a^{2}+s}-\frac{y^{2}}{b^{2}+s}
  -\frac{z^{2}}{c^{2}+s}\right)\frac{\mathrm{d}s}{R_s},\\
  &R_s=\sqrt{(a^{2}+s)(b^{2}+s)(c^{2}+s)},
\end{aligned}
\end{equation}
式中 $\xi$ 是方程
\begin{equation}
  \frac{x^{2}}{a^{2}+\xi}+\frac{y^{2}}{b^{2}+\xi}
  +\frac{z^{2}}{c^{2}+\xi}=1
\end{equation}

的正根.椭球内部场的势由如下公式给出：
\begin{equation}
  \varphi=-\pi\mu abck\int_{0}^{\infty}
  \left(1-\frac{x^{2}}{a^{2}+s}-\frac{y^{2}}{b^{2}+s}
  -\frac{z^{2}}{c^{2}+s}\right)\frac{\mathrm{d}s}{R_s},
\end{equation}

它同 (99.11) 的区别在于积分下限换成了 0；我们注意到这个表达式是坐标 $x,y,z$ 的二次函数.

按照 (99.6)，物体的引力能可通过将表达式 (99.13) 对椭球体积进行积分求得.这个积分可以用初等方法完成\footnote{对平方 $x^{2},y^{2},z^{2}$ 积分最简单的做法是，通过代换 $x=ax',\ y=by',\ z=cz'$ 把对椭球体积的积分化为对单位球体积的积分.}，结果是：
\begin{equation*}
\begin{aligned}
  U&=\frac{3km^{2}}{8}\int_{0}^{\infty}
  \left[\frac{1}{5}\left(\frac{a^{2}}{a^{2}+s}
  +\frac{b^{2}}{b^{2}+s}+\frac{c^{2}}{c^{2}+s}\right)-1\right]
  \frac{\mathrm{d}s}{R_s}\\
  &=\frac{3km^{2}}{8}\int_{0}^{\infty}
  \left[\frac{2}{5}s\,\mathrm{d}\left(\frac{1}{R_s}\right)
  -\frac{2}{5}\frac{\mathrm{d}s}{R_s}\right]
\end{aligned}
\end{equation*}
（$m=\dfrac{4\pi}{3}abc\,\mu$ 是物体的总质量）；对第一项作分部积分，最后得到：
\begin{equation}
  U=-\frac{3km^{2}}{10}\int_{0}^{\infty}\frac{\mathrm{d}s}{R_s}.
\end{equation}

出现在 (99.11)—(99.14) 式中的所有积分，都可以用第一类和第二类椭圆积分来表示.对于旋转椭球体，这些积分可以用初等函数来表示.就特例言之，旋转扁椭球（$a=b>c$）的引力能是
\begin{equation}
  U=-\frac{3km^{2}}{5\sqrt{a^{2}-c^{2}}}
  \arccos\frac{c}{a},
\end{equation}

而对于旋转长椭球（$a>b=c$）有：
\begin{equation}
  U=-\frac{3km^{2}}{5\sqrt{a^{2}-c^{2}}}
  \operatorname{arch}\frac{a}{c}.
\end{equation}

对于球（$a=c$），两个公式都给出值 $U=-3km^{2}/5a$，这个结果当然也可由初等方法得到\footnote{半径为 $a$ 的均匀球内部场的势为 $\varphi=-2\pi k\mu\left(a^{2}-\dfrac{r^{2}}{3}\right)$.}.

\begin{xiti}

\noindent{\heiti 习题}\quad 求作整体旋转、且引力质量分布均匀的流体的平衡形状.

解：平衡条件是，物体表面上引力势与离心力势之和为常量：
\begin{equation*}
  \varphi-\frac{\Omega^{2}}{2}(x^{2}+y^{2})=\mathrm{const}
\end{equation*}
（$\Omega$ 是角速度；旋转轴是 $z$ 轴）.满足要求的形状是旋转扁椭球.为求得它的参数，将 (99.13) 代入平衡条件，并用方程 (99.10) 消去 $z^{2}$；这就给出：
\begin{equation*}
\begin{aligned}
  &(x^{2}+y^{2})\left[\int_{0}^{\infty}
  \frac{\mathrm{d}s}{(a^{2}+s)^{2}\sqrt{c^{2}+s}}
  -\frac{\Omega^{2}}{2\pi\mu ka^{2}c}
  -\frac{c^{2}}{a^{2}}\int_{0}^{\infty}
  \frac{\mathrm{d}s}{(a^{2}+s)(c^{2}+s)^{3/2}}\right]\\
  &={}\mathrm{const},
\end{aligned}
\end{equation*}

由此得出，方括号内的表达式必须为零.进行积分，结果得到方程
\begin{equation*}
  \frac{a^{2}+2c^{2}}{(a^{2}-c^{2})^{3/2}}\arccos\frac{c}{a}
  -\frac{3c^{2}}{a^{2}-c^{2}}
  =\frac{\Omega^{2}}{2\pi k\mu}
  =\frac{25}{6}\left(\frac{4\pi}{3}\right)^{1/3}
  \frac{M^{2}\mu^{1/3}}{m^{10/3}k}\left(\frac{c}{a}\right)^{4/3}
\end{equation*}
（$M=\dfrac{2}{5}ma^{2}\Omega$ 是物体绕 $z$ 轴的角动量），该方程对于给定的 $M$ 或 $\Omega$ 决定了两半轴之比 $c/a$.比值 $c/a$ 对 $M$ 的依赖关系是单值的；$c/a$ 随着 $M$ 的增加而单调增加.

不过，我们得到的对称形状仅对不太大的 $M$ 值（对于小扰动）是稳定的.对于 $M=0.24\allowbreak\,k^{1/2}m^{5/3}\mu^{-1/6}$（$c/a=0.58$）将失去稳定性.随着 $M$ 的进一步增加，平衡形状变为 $b/a$ 和 $c/a$ 值（分别从 1 和 0.58）逐渐减小的一般椭球.这种形状对于 $M=0.31\allowbreak\,k^{1/2}m^{5/3}\mu^{-1/6}$（$a:b:c=1:0.43:0.34$）再次变得不稳定\footnote{有关这个问题的参考文献可在 H. Lamb 的著作 Hydrodynamics, XII, 1947 中找到.}.

\end{xiti}

\section{中心对称的引力场}

我们来研究一个具有中心对称的引力场.任何中心对称分布的物质都可以产生这样的场；这时，当然不仅分布必须是中心对称的，而且物质的运动也必须是中心对称的，就是说，每一点的速度必须是沿着径向的.

场的中心对称性就等于说，时空的度规，即间隔 $\mathrm{d}s$ 的表达式，在所有与中心等距离之点必须一样.在欧氏空间中，这个距离等于径矢；在非欧氏空间中，例如在引力场存在的情形下，没有一个量具有欧氏径矢的一切特性（例如，既等于到中心的距离，又等于圆周之长除以 $2\pi$）.因此，“径矢”的选择现在是任意的.

假如我们用空间“球”坐标 $r,\theta,\varphi$，那么，$\mathrm{d}s^{2}$ 的最广义的中心对称式是
\begin{equation}
\begin{aligned}
  \mathrm{d}s^{2}&=h(r,t)\,\mathrm{d}r^{2}
  +k(r,t)(\sin^{2}\theta\,\mathrm{d}\varphi^{2}+\mathrm{d}\theta^{2})\\
  &\quad+l(r,t)\,\mathrm{d}t^{2}+a(r,t)\,\mathrm{d}r\,\mathrm{d}t,
\end{aligned}
\end{equation}
式中，$a,h,k,l$ 是“径矢” $r$ 和“时间” $t$ 的某种函数.但是，因为在广义相对论中，参考系是可以任意选择的，所以我们还可以对坐标进行不破坏 $\mathrm{d}s^{2}$ 的中心对称性的任何变换；这就是说，我们可以按照公式
\begin{equation*}
  r=f_1(r',t'),\qquad t=f_2(r',t'),
\end{equation*}
变换坐标 $r$ 和 $t$，此处的 $f_1,f_2$ 是新坐标 $r',t'$ 的任意函数.

利用这种可能性，我们这样来选择坐标 $r$ 和时间 $t$：首先，使 $\mathrm{d}s^{2}$ 的表达式中的 $\mathrm{d}r\,\mathrm{d}t$ 的系数 $a(r,t)$ 为零，其次，使 $\mathrm{d}s^{2}$ 的表达式中的第二项的系数 $k(r,t)$ 简化为 $-r^{2}$.\footnote{必须指出这个条件并未唯一决定时间坐标的选择.就是说，它还可以进行一个不包含 $r$ 的形如 $t=f(t')$ 的任意变换.}后一要求意味着，径矢 $r$ 是这样定义的：使一个以坐标原点为圆心的圆周之长等于 $2\pi r$（在 $\theta=\pi/2$ 的平面内的圆的弧元等于 $\mathrm{d}l=r\,\mathrm{d}\varphi$）.将 $h$ 和 $l$ 分别写成形如 $-\mathrm{e}^{\lambda}$ 和 $c^{2}\mathrm{e}^{\nu}$ 的指数式是便利的，此处的 $\lambda$ 和 $\nu$ 是 $r$ 和 $t$ 的某种函数.因此，我们得到下面的 $\mathrm{d}s^{2}$ 表达式：
\begin{equation}
  \mathrm{d}s^{2}=\mathrm{e}^{\nu}c^{2}\mathrm{d}t^{2}
  -r^{2}(\mathrm{d}\theta^{2}+\sin^{2}\theta\,\mathrm{d}\varphi^{2})
  -\mathrm{e}^{\lambda}\mathrm{d}r^{2}.
\end{equation}

用 $x^0,x^1,x^2,x^3$ 分别代表坐标 $ct,r,\theta,\varphi$，对于度规张量的非零的分量，我们有以下表达式：
\begin{equation*}
  g_{00}=\mathrm{e}^{\nu},\quad g_{11}=-\mathrm{e}^{\lambda},\quad
  g_{22}=-r^{2},\quad g_{33}=-r^{2}\sin^{2}\theta.
\end{equation*}

显然
\begin{equation*}
  g^{00}=\mathrm{e}^{-\nu},\quad g^{11}=-\mathrm{e}^{-\lambda},\quad
  g^{22}=-r^{-2},\quad g^{33}=-r^{-2}\sin^{-2}\theta.
\end{equation*}

用这些值，不难从公式 (86.3) 算出 $\Gamma_{kl}^i$.计算的结果如下（一撇表示对于 $r$ 微分，在符号上的一点表示对于 $ct$ 微分）：
\begin{equation}
\begin{aligned}
  &\Gamma_{11}^{1}=\frac{\lambda'}{2},\quad
  \Gamma_{10}^{0}=\frac{\nu'}{2},\quad
  \Gamma_{33}^{2}=-\sin\theta\cos\theta,\\
  &\Gamma_{11}^{0}=\frac{\dot{\lambda}}{2}\mathrm{e}^{\lambda-\nu},\quad
  \Gamma_{22}^{1}=-r\mathrm{e}^{-\lambda},\quad
  \Gamma_{00}^{1}=\frac{\nu'}{2}\mathrm{e}^{\nu-\lambda},\\
  &\Gamma_{12}^{2}=\Gamma_{13}^{3}=\frac{1}{r},\quad
  \Gamma_{23}^{3}=\cot\theta,\quad
  \Gamma_{00}^{0}=\frac{\dot{\nu}}{2},\\
  &\Gamma_{10}^{1}=\frac{\dot{\lambda}}{2},\quad
  \Gamma_{33}^{1}=-r\sin^{2}\theta\,\mathrm{e}^{-\lambda}.
\end{aligned}
\end{equation}

$\Gamma_{kl}^i$ 的所有其余的分量（除了那些只需交换指标 $k$ 和 $l$ 就能从所写出的分量得到者外）都为零.

为了得到引力的方程，我们应当按照公式 (92.7) 计算张量 $R_k^i$ 的分量.简单的计算导出下面的方程：
\begin{equation}
  \frac{8\pi k}{c^{4}}T_1^1
  =-\mathrm{e}^{-\lambda}\left(\frac{\nu'}{r}+\frac{1}{r^{2}}\right)
  +\frac{1}{r^{2}},
\end{equation}

\begin{equation}
\begin{aligned}
  \frac{8\pi k}{c^{4}}T_2^2=\frac{8\pi k}{c^{4}}T_3^3
  &=-\frac{1}{2}\mathrm{e}^{-\lambda}
  \left(\nu''+\frac{\nu'^{2}}{2}
  +\frac{\nu'-\lambda'}{r}-\frac{\nu'\lambda'}{2}\right)\\
  &\quad+\frac{1}{2}\mathrm{e}^{-\nu}
  \left(\ddot{\lambda}+\frac{\dot{\lambda}^{2}}{2}
  -\frac{\dot{\lambda}\dot{\nu}}{2}\right),
\end{aligned}
\end{equation}

\begin{equation}
  \frac{8\pi k}{c^{4}}T_0^0
  =-\mathrm{e}^{-\lambda}\left(\frac{1}{r^{2}}-\frac{\lambda'}{r}\right)
  +\frac{1}{r^{2}},
\end{equation}

\begin{equation}
  \frac{8\pi k}{c^{4}}T_0^1=-\mathrm{e}^{-\lambda}\frac{\dot{\lambda}}{r}
\end{equation}
（(95.6) 的其余分量恒等于零）.利用 (94.9)，能量动量张量可以由物质的能量密度 $\varepsilon$，它的压强 $p$，和径向速度 $v$ 来表示.

在真空中，即在产生场的质量以外，方程 (100.4)—(100.7) 对于中心对称场这一非常重要的情形是严格可积的.令能量动量张量等于零，我们得到下面的方程：
\begin{equation}
  \mathrm{e}^{-\lambda}\left(\frac{\nu'}{r}+\frac{1}{r^{2}}\right)
  -\frac{1}{r^{2}}=0,
\end{equation}

\begin{equation}
  \mathrm{e}^{-\lambda}\left(\frac{\lambda'}{r}-\frac{1}{r^{2}}\right)
  +\frac{1}{r^{2}}=0,
\end{equation}

\begin{equation}
  \dot{\lambda}=0
\end{equation}
（我们没有写第四个方程，即方程 (100.5)，因为它可以从其余三个方程推出）.

从 (100.10) 式直接看出，$\lambda$ 与时间无关.此处，将 (100.8) 和 (100.9) 相加，可得 $\lambda'+\nu'=0$，即
\begin{equation}
  \lambda+\nu=f(t),
\end{equation}
式中，$f(t)$ 仅仅是时间的函数.但是当我们选择间隔 $\mathrm{d}s^{2}$ 为 (100.2) 的形式后，仍旧可以作一个形如 $t=f(t')$ 的任意的时间变换.这样的变换就等效于将一个任意的时间函数加到 $\nu$ 上，并且借它的帮助，我们总可能使 (100.11) 式中 $f(t)$ 为零.因此，不失任何普遍性，我们可以认为 $\lambda+\nu=0$.由此我们注意到，真空中的中心对称引力场自然而然是静态的.

方程 (100.9) 是容易积分的，其结果是
\begin{equation}
  \mathrm{e}^{-\lambda}=\mathrm{e}^{\nu}
  =1+\frac{\mathrm{const}}{r}.
\end{equation}

因此，在无限远处（$r\to\infty$）$\mathrm{e}^{-\lambda}=\mathrm{e}^{\nu}=1$，这就是说，在与引力物体相距很远之处，度规就自然而然地变成了伽利略度规.上式中的常数容易用物体的质量来表示，只要要求在远距离处，场是弱的，牛顿定律应当成立.\footnote{对于中心对称分布的球形空腔内部的场，必有 $\mathrm{const}=0$，因为若不如此，度规在 $r=0$ 处会有奇异性.因此，这样的空腔内部的度规自然是伽利略度规，即空腔内部没有引力场（正如在牛顿理论中一样）.}换言之，我们应当有 $g_{00}=1+2\varphi/c^{2}$，此处的势 $\varphi$ 具有牛顿力学中的值 (99.4)：$\varphi=-km/r$（$m$ 是产生场的物体的总质量）.从此显然可见，(100.12) 式中的 $\mathrm{const}=-2km/c^{2}$.这个量具有长度的量纲，称为物体的引力半径 $r_{\mathrm{g}}$：
\begin{equation}
  r_{\mathrm{g}}=\frac{2km}{c^{2}}.
\end{equation}

因此，我们最后得到如下形式的时空度规：
\begin{equation}
  \mathrm{d}s^{2}=\left(1-\frac{r_{\mathrm{g}}}{r}\right)c^{2}\mathrm{d}t^{2}
  -r^{2}(\sin^{2}\theta\,\mathrm{d}\varphi^{2}+\mathrm{d}\theta^{2})
  -\frac{\mathrm{d}r^{2}}{1-r_{\mathrm{g}}/r}.
\end{equation}

爱因斯坦方程的这个解是史瓦西得到的 (K. Schwarzschild, 1916).它完全决定了任何中心对称分布的物质在真空中产生的引力场.我们着重指出，这个解不仅对于静止的物质有效，而且对于运动的物质也有效，只要求这个运动也有中心对称性（例如，中心对称的脉动）.我们注意到，度规 (100.14) 只依赖于引力物体的总质量，恰如牛顿理论中类似的问题.

空间度规由空间距离元的表达式
\begin{equation}
  \mathrm{d}l^{2}=\frac{\mathrm{d}r^{2}}{1-r_{\mathrm{g}}/r}
  +r^{2}(\sin^{2}\theta\,\mathrm{d}\varphi^{2}+\mathrm{d}\theta^{2})
\end{equation}
所决定，坐标 $r$ 的几何意义决定于如下事实，即在度规 (100.15) 中，其中心在场中心的圆的周长是 $2\pi r$.但同一半径上 $r_1$ 和 $r_2$ 两点之间的距离由以下积分给出：
\begin{equation}
  \int_{r_1}^{r_2}\frac{\mathrm{d}r}{\sqrt{1-r_{\mathrm{g}}/r}}
  >r_2-r_1.
\end{equation}

此外，我们看到，$g_{00}\leqslant1$.结合定义固有时的公式 (84.1) $\mathrm{d}\tau=\sqrt{g_{00}}\,\mathrm{d}t$ 可以推断：
\begin{equation}
  \mathrm{d}\tau\leqslant\mathrm{d}t.
\end{equation}

只有在无穷远处，等号方能成立，在那里，$t$ 与固有时重合.因此，在与物质相距有限距离处的时间比在无穷远处的时间“走得慢些”.

最后我们再介绍 $\mathrm{d}s^{2}$ 在距离坐标原点很远处的一个近似式：
\begin{equation}
  \mathrm{d}s^{2}=\mathrm{d}s_0^{2}
  -\frac{2km}{c^{2}r}(\mathrm{d}r^{2}+c^{2}\mathrm{d}t^{2}).
\end{equation}

第二项是对伽利略度规 $\mathrm{d}s_0^{2}$ 的一个小的修正.在与产生场的物质相距甚远处，每个场都是中心对称的.因此，(100.18) 式决定与任何物体体系相距甚远之处的度规.

在讨论引力物质内部的中心对称引力场时，也可以作某些一般的考虑.从方程 (100.6)，我们看出，当 $r\to0$ 时，$\lambda$ 至少要像 $r^{2}$ 一样快地趋近于零；假如不是如此，那么方程右边当 $r\to0$ 时就应该变为无穷大了，就是说，$T_0^0$ 在 $r=0$ 处应该有一个奇点，然而这在物理上是荒谬的.将 (100.6) 作形式积分，其边界条件是 $\lambda|_{r=0}=0$，我们得到
\begin{equation}
  \lambda=-\ln\left(1-\frac{8\pi k}{c^{4}r}
  \int_{0}^{r}T_0^0r^{2}\,\mathrm{d}r\right).
\end{equation}

因从 (94.10)，有 $T_0^0=\mathrm{e}^{-\nu}T_{00}\geqslant0$，由此可见 $\lambda\geqslant0$，即
\begin{equation}
  \mathrm{e}^{\lambda}\geqslant1.
\end{equation}

从 (100.4) 中逐项减去方程 (100.6)，我们得到：
\begin{equation*}
  \frac{\mathrm{e}^{-\lambda}}{r}(\nu'+\lambda')
  =\frac{8\pi k}{c^{4}}(T_0^0-T_1^1)
  =\frac{\displaystyle(\varepsilon+p)\left(1+\frac{v^{2}}{c^{2}}\right)}
  {\displaystyle1-\frac{v^{2}}{c^{2}}}\geqslant0,
\end{equation*}
即 $\nu'+\lambda'\geqslant0$.但对于 $r\to\infty$（远离物质），度规变为伽利略度规，即 $\nu\to0$，$\lambda\to0$.因此，从 $\nu'+\lambda'\geqslant0$ 可以断定，在整个空间中都有
\begin{equation}
  \nu+\lambda\leqslant0.
\end{equation}

因为 $\lambda\geqslant0$，从而得到 $\nu\leqslant0$，即
\begin{equation}
  \mathrm{e}^{\nu}\leqslant1.
\end{equation}

这些不等式表明，前面所讲的真空中中心对称场的空间度规的性质 (100.16) 和 (100.17) 以及钟的快慢同样可以应用于引力物质内部的场.

如果引力场是由一个“半径”为 $a$ 的球体所产生，那么，当 $r>a$ 时，我们就有 $T_0^0=0$.对于 $r>a$ 的点，公式 (100.19) 给出
\begin{equation*}
  \lambda=-\ln\left(1-\frac{8\pi k}{c^{4}r}
  \int_{0}^{a}T_0^0r^{2}\,\mathrm{d}r\right).
\end{equation*}

另一方面，在这里我们可以用适于真空的表达式 (100.14)，按照此式，
\begin{equation*}
  \lambda=-\ln\left(1-\frac{2km}{c^{2}r}\right).
\end{equation*}

使以上两式相等，我们得到公式
\begin{equation}
  m=\frac{4\pi}{c^{2}}\int_{0}^{a}T_0^0r^{2}\,\mathrm{d}r,
\end{equation}

这个公式是用物体的能量动量张量来表示物体的总质量.

就特例而言，对于物体内物质的静态分布，我们有 $T_0^0=\varepsilon$，因此有
\begin{equation}
  m=\frac{4\pi}{c^{4}}\int_{0}^{a}\varepsilon r^{2}\,\mathrm{d}r.
\end{equation}

请注意，积分是对 $4\pi r^{2}\mathrm{d}r$ 进行的，而度规 (100.2) 的空间体积元是
\begin{equation*}
  \mathrm{d}V=4\pi r^{2}\mathrm{e}^{\lambda/2}\,\mathrm{d}r,
\end{equation*}
这里，按照 (100.20)，$\mathrm{e}^{\lambda/2}>1$.这个差表示了物体的引力质量亏损.

\begin{xiti}

\noindent{\heiti 习题1}\quad 求史瓦西度规 (100.14) 的曲率张量的不变量.

解：用 (92.1) 式和来自 (100.3) 的 $\Gamma_{kl}^i$（或用§92习题2获得的公式）进行计算，得到曲率张量非零分量的下列值：
\begin{equation*}
  R_{0101}=\frac{r_{\mathrm{g}}}{r^{3}},\qquad
  R_{0202}=\frac{R_{0303}}{\sin^{2}\theta}
  =-\frac{r_{\mathrm{g}}(r-r_{\mathrm{g}})}{2r^{2}},
\end{equation*}

\begin{equation*}
  R_{1212}=\frac{R_{1313}}{\sin^{2}\theta}
  =\frac{r_{\mathrm{g}}}{2(r-r_{\mathrm{g}})},\qquad
  R_{2323}=-r\,r_{\mathrm{g}}\sin^{2}\theta.
\end{equation*}

对于 (92.20) 的不变量 $I_1$ 和 $I_2$，我们得到：
\begin{equation*}
  I_1=\left(\frac{r_{\mathrm{g}}}{2r^{3}}\right)^{2},\qquad
  I_2=-\left(\frac{r_{\mathrm{g}}}{2r^{3}}\right)^{3}
\end{equation*}
（包含对偶张量 $\stackrel{*}{R}_{iklm}$ 的乘积恒等于零）.该曲率张量属于彼得罗夫分类的 D 型（有实不变量 $\lambda^{(1)}=\lambda^{(2)}=-r_{\mathrm{g}}/2r^{3}$）.我们注意到，该曲率不变量仅在 $r=0$，而不是在 $r=r_{\mathrm{g}}$ 有奇异性.

\noindent{\heiti 习题2}\quad 求这个度规的空间曲率.

解：空间曲率张量 $P_{\alpha\beta\gamma\delta}$ 的分量可用张量 $P_{\alpha\beta}$（和张量 $\gamma_{\alpha\beta}$）的分量来表示.因此，只须计算 $P_{\alpha\beta}$ 就够了（参见§92习题1）.利用 $\gamma_{\alpha\beta}$ 来表示张量 $P_{\alpha\beta}$，正如用 $g_{ik}$ 来表示 $R_{ik}$ 一样.用从 (100.15) 得到的 $\gamma_{\alpha\beta}$ 值，经过计算后，我们得到：
\begin{equation*}
  P_{\theta}^{\theta}=P_{\varphi}^{\varphi}
  =\frac{r_{\mathrm{g}}}{2r^{3}},\qquad
  P_{r}^{r}=-\frac{r_{\mathrm{g}}}{r^{3}},
\end{equation*}
当 $\alpha\neq\beta$ 时，$P_{\alpha}^{\beta}=0$.我们注意到 $P_{\theta}^{\theta},P_{\varphi}^{\varphi}>0$，$P_r^r<0$，而 $P\equiv P_{\alpha}^{\alpha}=0$.

从§92习题1给出的公式，我们得到：
\begin{equation*}
\begin{aligned}
  &P_{r\theta r\theta}=(P_r^r+P_\theta^\theta)\gamma_{rr}\gamma_{\theta\theta}
  =-P_{\varphi}^{\varphi}\gamma_{rr}\gamma_{\theta\theta},\\
  &P_{r\varphi r\varphi}=-P_{\theta}^{\theta}\gamma_{rr}\gamma_{\varphi\varphi},
  \qquad
  P_{\theta\varphi\theta\varphi}
  =-P_r^r\gamma_{\theta\theta}\gamma_{\varphi\varphi}.
\end{aligned}
\end{equation*}

于是推得（参见§92第四个脚注），对垂直于半径的“平面”，高斯曲率是
\begin{equation*}
  K=\frac{P_{\theta\varphi\theta\varphi}}
  {\gamma_{\theta\theta}\gamma_{\varphi\varphi}}
  =-P_r^r>0
\end{equation*}
（这意味着，在该“平面”上与垂直于它的半径的交点的邻域内画一个小三角形，该三角形的三个角之和大于 $\pi$）.至于通过中心的“平面”，它们的高斯曲率 $K<0$，这意味着，这样一个“平面”内的小三角形的三个角之和小于 $\pi$（然而，这并不是指围绕中心的三角形，这样的三角形三个角之和大于 $\pi$）.

\noindent{\heiti 习题3}\quad 决定旋转曲面的形式，我们要求在这个旋转曲面上的几何同在真空中的中心对称引力场内经过原点的“平面”上的几何一样.

解：在旋转曲面 $z=z(r)$ 上的几何为（用柱体坐标）线元
\begin{equation*}
  \mathrm{d}l^{2}=\mathrm{d}r^{2}+\mathrm{d}z^{2}+r^{2}\mathrm{d}\varphi^{2}
  =\mathrm{d}r^{2}(1+z'^{2})+r^{2}\mathrm{d}\varphi^{2}
\end{equation*}
所决定.将此式与所考虑的场内的“平面” $\theta=\pi/2$ 内的线元 (100.15)
\begin{equation*}
  \mathrm{d}l^{2}=r^{2}\mathrm{d}\varphi^{2}
  +\frac{\mathrm{d}r^{2}}{1-r_{\mathrm{g}}/r}
\end{equation*}
相比较，我们得到
\begin{equation*}
  1+z'^{2}=\left(1-\frac{r_{\mathrm{g}}}{r}\right)^{-1},
\end{equation*}
由此可得
\begin{equation*}
  z=2\sqrt{r_{\mathrm{g}}(r-r_{\mathrm{g}})}.
\end{equation*}

对于 $r=r_{\mathrm{g}}$，这个函数有一个奇异性——分支点.其解释是，空间度规 (100.15) 与时空度规 (100.14) 不同，在 $r=r_{\mathrm{g}}$ 确实有一个奇异性.

上题提到的经过中心的“平面”上的一般几何性质，也可以通过考虑这里给出的直观模型中的曲率得到.

\noindent{\heiti 习题4}\quad 将间隔 (100.14) 变换到这样的坐标中，在其中的空间距离元具有共形欧几里得形式（即 $\mathrm{d}l^{2}$ 和它的欧氏表达式成比例）.

解：令
\begin{equation*}
  r=\rho\left(1+\frac{r_{\mathrm{g}}}{4\rho}\right)^{2},
\end{equation*}
从 (100.14)，我们得到
\begin{equation*}
  \mathrm{d}s^{2}=\left(\frac{1-\dfrac{r_{\mathrm{g}}}{4\rho}}
  {1+\dfrac{r_{\mathrm{g}}}{4\rho}}\right)^{2}c^{2}\mathrm{d}t^{2}
  -\left(1+\frac{r_{\mathrm{g}}}{4\rho}\right)^{4}
  (\mathrm{d}\rho^{2}+\rho^{2}\mathrm{d}\theta^{2}
  +\rho^{2}\sin^{2}\theta\,\mathrm{d}\varphi^{2}).
\end{equation*}

坐标 $\rho,\theta,\varphi$ 称为各向同性球坐标.我们也可以引入各向同性笛卡儿坐标 $x,y,z$ 来代替它们.特别是，在大距离处（$\rho\gg r_{\mathrm{g}}$），我们近似地有
\begin{equation*}
  \mathrm{d}s^{2}=\left(1-\frac{r_{\mathrm{g}}}{\rho}\right)c^{2}\mathrm{d}t^{2}
  -\left(1+\frac{r_{\mathrm{g}}}{\rho}\right)
  (\mathrm{d}x^{2}+\mathrm{d}y^{2}+\mathrm{d}z^{2}).
\end{equation*}

\noindent{\heiti 习题5}\quad 求在共动参考系中物质内部的中心对称引力场方程.

解：我们利用间隔元 (100.1) 内的坐标 $r,t$ 的两个可能的变换，使得：第一，使 $\mathrm{d}r\,\mathrm{d}t$ 的系数 $a(r,t)$ 为零；第二，使物体的径向速度在所有的点为零（根据中心对称性，速度的其余分量一般都为零）.在此以后，坐标 $r$ 和 $t$ 仍然可能经受一个形如 $r=r(r'),\ t=t(t')$ 的任意变换.

将这样选择的径向坐标和时间记为 $R$ 和 $\tau$，分别用 $-\mathrm{e}^{\lambda},-\mathrm{e}^{\mu},\mathrm{e}^{\nu}$（$\lambda,\mu,\nu$ 是 $R$ 和 $\tau$ 的函数）来表示 $h,k,l$，我们得到线元的表达式如下：
\begin{equation}
  \mathrm{d}s^{2}=c^{2}\mathrm{e}^{\nu}\mathrm{d}\tau^{2}
  -\mathrm{e}^{\lambda}\mathrm{d}R^{2}
  -\mathrm{e}^{\mu}(\mathrm{d}\theta^{2}+\sin^{2}\theta\,\mathrm{d}\varphi^{2})
  \tag{1}
\end{equation}

在共动参考系中物质的能量动量张量的分量是：
\begin{equation*}
  T_0^0=\varepsilon,\qquad T_1^1=T_2^2=T_3^3=-p.
\end{equation*}

计算给出下面的场方程\footnote{分量 $R_{ik}$ 可以按正文中做过的那样直接计算，或者用§92习题2得到的公式计算.}：
  \begin{equation}
  -\frac{8\pi k}{c^{4}}T_1^1=\frac{8\pi k}{c^{4}}p
  =\frac{1}{2}\mathrm{e}^{-\lambda}
  \left(\frac{\mu'^{2}}{2}+\mu'\nu'\right)
  -\mathrm{e}^{-\nu}\left(\ddot{\mu}-\frac{1}{2}\dot{\mu}\dot{\nu}
  +\frac{3}{4}\dot{\mu}^{2}\right)-\mathrm{e}^{-\mu},
  \tag{2}
\end{equation}

\begin{equation}
\begin{aligned}
  -\frac{8\pi k}{c^{4}}T_2^2=\frac{8\pi k}{c^{4}}p
  ={}&\frac{1}{4}\mathrm{e}^{-\lambda}
  (2\nu''+\nu'^{2}+2\mu''+\mu'^{2}
  -\mu'\lambda'-\nu'\lambda'+\mu'\nu')\\
  &+\frac{1}{4}\mathrm{e}^{-\nu}
  (\dot{\lambda}\dot{\nu}+\dot{\mu}\dot{\nu}
  -\dot{\lambda}\dot{\mu}-2\ddot{\lambda}-\dot{\lambda}^{2}
  -2\ddot{\mu}-\dot{\mu}^{2}),
\end{aligned}
\tag{3}
\end{equation}

\begin{equation}
  \frac{8\pi k}{c^{4}}T_0^0=\frac{8\pi k}{c^{4}}\varepsilon
  =-\mathrm{e}^{-\lambda}\left(\mu''+\frac{3}{4}\mu'^{2}
  -\frac{\mu'\lambda'}{2}\right)
  +\frac{1}{2}\mathrm{e}^{-\nu}\left(\dot{\lambda}\dot{\mu}
  +\frac{\dot{\mu}^{2}}{2}\right)+\mathrm{e}^{-\mu},
  \tag{4}
\end{equation}

\begin{equation}
  \frac{8\pi k}{c^{4}}T_0^1=0
  =\frac{1}{2}\mathrm{e}^{-\lambda}
  (2\dot{\mu}'+\dot{\mu}\mu'-\dot{\lambda}\mu'-\nu'\dot{\mu})
  \tag{5}
\end{equation}
（式中撇号表示对 $R$ 微分，点表示对 $c\tau$ 微分）.

从包含在引力场方程之内的方程 $T_{i;k}^k=0$ 出发，不难得到 $\lambda,\mu,\nu$ 各量间的一些普遍关系.用公式 (86.11)，我们得到下面两个方程：
\begin{equation}
  \dot{\lambda}+2\dot{\mu}=-\frac{2\dot{\varepsilon}}{p+\varepsilon},
  \qquad
  \nu'=-\frac{2p'}{p+\varepsilon}.
  \tag{6}
\end{equation}

如果 $p$ 是能量 $\varepsilon$ 的已知函数，那么，方程 (6) 可以直接积分如：
\begin{equation}
  \lambda+2\mu=-2\int\frac{\mathrm{d}\varepsilon}{p+\varepsilon}
  +f_1(R),\qquad
  \nu=-2\int\frac{\mathrm{d}p}{p+\varepsilon}+f_2(\tau),
  \tag{7}
\end{equation}
鉴于上面所说的可以作形如 $R=R(R'),\ \tau=\tau(\tau')$ 的任意变换，上式中的函数 $f_1(R)$ 和 $f_2(\tau)$ 可以任意选择.

\noindent{\heiti 习题6}\quad 求决定静止轴对称物体周围真空中的静态引力场的方程 (H. Weyl, 1917).

解：假设空间柱坐标 $x^1=\varphi,\ x^2=\rho,\ x^3=z$ 中的静态线元具有形式
\begin{equation*}
  \mathrm{d}s^{2}=\mathrm{e}^{\nu}c^{2}\mathrm{d}t^{2}
  -\mathrm{e}^{\omega}\mathrm{d}\varphi^{2}
  -\mathrm{e}^{\mu}(\mathrm{d}\rho^{2}+\mathrm{d}z^{2}),
\end{equation*}
式中 $\nu,\omega,\mu$ 是 $\rho$ 和 $z$ 的函数；这样的表示把坐标的选择确定到形如 $\rho=\rho(\rho',z'),\ z=z(\rho',z')$ 的变换之内，它仅给二次型 $\mathrm{d}\rho^{2}+\mathrm{d}z^{2}$ 乘上一个共同的因子.

从方程
\begin{equation*}
  R_0^0=\frac{1}{4}\mathrm{e}^{-\mu}
  \left[2\nu_{,\rho,\rho}+\nu_{,\rho}(\nu_{,\rho}+\omega_{,\rho})
  +2\nu_{,z,z}+\nu_{,z}(\nu_{,z}+\omega_{,z})\right]=0,
\end{equation*}

\begin{equation*}
  R_1^1=\frac{1}{4}\mathrm{e}^{-\mu}
  \left[2\omega_{,\rho,\rho}+\omega_{,\rho}(\nu_{,\rho}+\omega_{,\rho})
  +2\omega_{,z,z}+\omega_{,z}(\nu_{,z}+\omega_{,z})\right]=0
\end{equation*}
（式中下标 $,\rho$ 和 $,z$ 表示对 $\rho$ 和 $z$ 微分），取其和，我们得到：
\begin{equation*}
  \rho'_{,\rho,\rho}+\rho'_{,z,z}=0,
\end{equation*}
式中
\begin{equation*}
  \rho'(\rho,z)=\mathrm{e}^{(\nu+\omega)/2}.
\end{equation*}

因此 $\rho'(\rho,z)$ 是变量 $\rho$ 和 $z$ 的调和函数.按照这种函数的熟知性质，这意味着存在一个共轭调和函数 $z'(\rho,z)$，使得 $\rho'+\mathrm{i}z'=f(\rho+\mathrm{i}z)$，式中 $f$ 是复变量 $\rho+\mathrm{i}z$ 的解析函数.如果我们现在选择 $\rho',z'$ 为新坐标，由于变换 $\rho,z\to\rho',z'$ 的共形性，我们将有
\begin{equation*}
  \mathrm{e}^{\mu}(\mathrm{d}\rho^{2}+\mathrm{d}z^{2})
  =\mathrm{e}^{\mu'}(\mathrm{d}\rho'^{2}+\mathrm{d}z'^{2}),
\end{equation*}
式中 $\mu'(\rho',z')$ 是某个新函数.同时 $\mathrm{e}^{\omega}=\rho'^{2}\mathrm{e}^{-\nu}$；记 $\omega+\nu=\gamma$，去掉撇号，我们将 $\mathrm{d}s^{2}$ 写为如下形式：
\begin{equation}
  \mathrm{d}s^{2}=\mathrm{e}^{\nu}c^{2}\mathrm{d}t^{2}
  -\rho^{2}\mathrm{e}^{-\nu}\mathrm{d}\varphi^{2}
  -\mathrm{e}^{\gamma-\nu}(\mathrm{d}\rho^{2}+\mathrm{d}z^{2}).
  \tag{1}
\end{equation}

对这个度规构建方程 $R_0^0=0,\ R_3^3-R_2^2=0,\ R_2^3=0$，我们得到：
\begin{equation}
  \frac{1}{\rho}\frac{\partial}{\partial\rho}
  \left(\rho\frac{\partial\nu}{\partial\rho}\right)
  +\frac{\partial^{2}\nu}{\partial z^{2}}=0,
  \tag{2}
\end{equation}

\begin{equation}
  \frac{\partial\gamma}{\partial z}
  =\rho\frac{\partial\nu}{\partial\rho}\frac{\partial\nu}{\partial z},
  \qquad
  \frac{\partial\gamma}{\partial\rho}
  =\frac{\rho}{2}\left[\left(\frac{\partial\nu}{\partial\rho}\right)^{2}
  -\left(\frac{\partial\nu}{\partial z}\right)^{2}\right].
  \tag{3}
\end{equation}

我们注意到，(2) 具有柱坐标中拉普拉斯方程的形式（对于不依赖于 $\varphi$ 的函数）.如果解出这个方程，那么函数 $\gamma(\rho,z)$ 就完全由方程 (2)，(3) 决定了.在离产生场的物体很远处，函数 $\nu$ 和 $\gamma$ 应当趋于零.

\end{xiti}
\section{中心对称引力场中的运动}

我们现在来研究粒子在中心对称引力场中的运动.正如在每一个中心对称场内一样，运动发生在经过原点的一个“平面”上；我们选择这个平面为平面 $\theta=\pi/2$.

为了决定粒子的轨道，我们用哈密顿-雅可比方程：
\begin{equation*}
  g^{ik}\frac{\partial S}{\partial x^{i}}
  \frac{\partial S}{\partial x^{k}}-m^{2}c^{2}=0,
\end{equation*}
式中 $m$ 为粒子的质量（中心物体的质量记为 $m'$）.利用表达式 (100.14) 给出的 $g^{ik}$，我们求出下面的方程：
\begin{equation}
  \left(1-\frac{r_{\mathrm{g}}}{r}\right)^{-1}
  \left(\frac{\partial S}{c\,\partial t}\right)^{2}
  -\left(1-\frac{r_{\mathrm{g}}}{r}\right)
  \left(\frac{\partial S}{\partial r}\right)^{2}
  -\frac{1}{r^{2}}\left(\frac{\partial S}{\partial\varphi}\right)^{2}
  -m^{2}c^{2}=0,
\end{equation}
其中，$r_{\mathrm{g}}=2m'k/c^{2}$ 是中心物体的引力半径.

用解哈密顿-雅可比方程的一般法则，我们来寻求下面形式的 $S$：
\begin{equation}
  S=-\mathcal{E}_0t+M\varphi+S_r(r),
\end{equation}
式中，$\mathcal{E}_0$ 是常能量，$M$ 是角动量.将 (101.2) 代入 (101.1)，我们便求得导数 $\mathrm{d}S_r/\mathrm{d}r$，因此：
\begin{equation}
  S_r=\int\left[\frac{\mathcal{E}_0^{2}}{c^{2}}
  \left(1-\frac{r_{\mathrm{g}}}{r}\right)^{-2}
  -\left(m^{2}c^{2}+\frac{M^{2}}{r^{2}}\right)
  \left(1-\frac{r_{\mathrm{g}}}{r}\right)^{-1}\right]^{1/2}\mathrm{d}r.
\end{equation}

依赖关系 $r=r(t)$ 由方程 $\partial S/\partial\mathcal{E}_0=\mathrm{const}$ 给出（参见本教程第一卷§47），由此得到
\begin{equation}
  ct=\frac{\mathcal{E}_0}{mc^{2}}\int
  \frac{\mathrm{d}r}{\left(1-\dfrac{r_{\mathrm{g}}}{r}\right)
  \left[\left(\dfrac{\mathcal{E}_0}{mc^{2}}\right)^{2}
  -\left(1+\dfrac{M^{2}}{m^{2}c^{2}r^{2}}\right)
  \left(1-\dfrac{r_{\mathrm{g}}}{r}\right)\right]^{1/2}},
\end{equation}

轨道为方程 $\partial S/\partial M=\mathrm{const}$ 所决定，由此可得
\begin{equation}
  \varphi=\int\frac{M}{r^{2}}
  \left[\frac{\mathcal{E}_0^{2}}{c^{2}}
  -\left(m^{2}c^{2}+\frac{M^{2}}{r^{2}}\right)
  \left(1-\frac{r_{\mathrm{g}}}{r}\right)\right]^{-1/2}\mathrm{d}r.
\end{equation}

这个积分可以化为椭圆积分.

对于行星在太阳引力场内的运动，因为行星的速度比光速小很多，所以引力的相对论理论同牛顿理论相比，只能导出一个并不显著的修正.在轨道方程 (101.5) 内的被积式中，这相应于比值 $r_{\mathrm{g}}/r$ 很小，这里 $r_{\mathrm{g}}$ 是太阳的引力半径\footnote{对于太阳，$r_{\mathrm{g}}=3\,\mathrm{km}$；对于地球，$r_{\mathrm{g}}=0.9\,\mathrm{cm}$.}.

为了计算轨道的相对论修正，比较便利的是从作用量径向部分的表达式 (101.3) 出发，然后再对 $M$ 微分.变换积分变量，记
\begin{equation*}
  r(r-r_{\mathrm{g}})=r'^{2}
  \qquad\text{因而}\qquad
  r-r_{\mathrm{g}}\cong r',
\end{equation*}
其结果是平方根号内带 $M^{2}$ 的项取 $M^{2}/r'^{2}$ 的形式.在其他的项中按 $r_{\mathrm{g}}/r'$ 的幂级数展开，在要求的精度内得到：
\begin{equation}
\begin{aligned}
  S_r=\int\Biggl[&\left(2\mathcal{E}'m+\frac{\mathcal{E}'^{2}}{c^{2}}\right)
  +\frac{1}{r}(2m^{2}m'k+4\mathcal{E}'mr_{\mathrm{g}})\\
  &-\frac{1}{r^{2}}\left(M^{2}
  -\frac{3m^{2}c^{2}r_{\mathrm{g}}^{2}}{2}\right)\Biggr]^{1/2}\mathrm{d}r,
\end{aligned}
\end{equation}
这里为简便起见，我们略去了 $r'$ 上的撇号，并引入了非相对论能量 $\mathcal{E}'$（没有静能）.

根号下前两项的系数的修正项只有不是特别有趣的效应，即改变粒子能量和动量之间的关系及改变其牛顿轨道（椭圆）的参数.但 $1/r^{2}$ 项的系数中的改变却导致更根本的效应，即轨道近日点的系统（长期）移动.

因为轨道是由方程 $\varphi+\partial S_r/\partial M=\mathrm{const}$ 决定的，行星在其轨道上转一圈后角 $\varphi$ 的改变是
\begin{equation*}
  \Delta\varphi=-\frac{\partial}{\partial M}\Delta S_r,
\end{equation*}
式中 $\Delta S_r$ 是 $S_r$ 的相应改变.将 $S_r$ 展开为 $1/r^{2}$ 的系数的微小修正的幂级数，我们得到：
\begin{equation*}
  \Delta S_r=\Delta S_r^{(0)}
  -\frac{3m^{2}c^{2}r_{\mathrm{g}}^{2}}{4M}
  \frac{\partial\Delta S_r^{(0)}}{\partial M},
\end{equation*}
式中 $\Delta S_r^{(0)}$ 相应于在不移动的封闭椭圆中的运动.将这个关系对 $M$ 进行微分，并考虑到
\begin{equation*}
  -\frac{\partial}{\partial M}\Delta S_r^{(0)}=\Delta\varphi^{(0)}=2\pi,
\end{equation*}
我们得到：
\begin{equation*}
  \Delta\varphi=2\pi+\frac{3\pi m^{2}c^{2}r_{\mathrm{g}}^{2}}{2M^{2}}
  =2\pi+\frac{6\pi k^{2}m^{2}m'^{2}}{c^{2}M^{2}}.
\end{equation*}

第二项是要求的转一圈后牛顿椭圆的角位移 $\delta\varphi$，即轨道近日点的移动.借助公式
\begin{equation*}
  \frac{M^{2}}{km'm^{2}}=a(1-e^{2}),
\end{equation*}
可以将该移动用轨道半长轴的长度 $a$ 和偏心率 $e$ 来表示.我们得到\footnote{从（101.7）式求得的近日点移动数值，对于水星和地球分别等于每世纪 $43.0''$ 和 $3.8''$.}：
\begin{equation}
  \delta\varphi=\frac{6\pi km'}{c^{2}a(1-e^{2})}.
\end{equation}

下面我们来考虑光线在中心对称引力场内的路径.这个路径为程函方程 (87.9)
\begin{equation*}
  g^{ik}\frac{\partial\psi}{\partial x^{i}}
  \frac{\partial\psi}{\partial x^{k}}=0
\end{equation*}
所决定，这个方程就是 $m=0$ 的哈密顿-雅可比方程.因此可以通过在 (101.5) 中令 $m=0$ 立即求得光线的轨道；同时，必须用光的频率 $\omega_0=-\partial\psi/\partial t$ 来代替粒子的能量 $\mathcal{E}_0=-\partial S/\partial t$，再引入由 $\rho=cM/\omega_0$ 定义的常量 $\rho$ 来代替 $M$，我们得到：
\begin{equation}
  \varphi=\int\frac{\mathrm{d}r}
  {r^{2}\sqrt{\dfrac{1}{\rho^{2}}
  -\dfrac{1}{r^{2}}\left(1-\dfrac{r_{\mathrm{g}}}{r}\right)}}.
\end{equation}

如果忽略相对论修正（$r_{\mathrm{g}}\to0$），这个方程给出 $r=\rho/\cos\varphi$，即离原点距离为 $\rho$ 的直线.为了研究相对论修正，我们以前例同样的方式进行.

对于程函的径向部分，我们有（参见 (101.3)）：
\begin{equation*}
  \psi_r(r)=\frac{\omega_0}{c}\int
  \sqrt{\frac{r^{2}}{(r-r_{\mathrm{g}})^{2}}
  -\frac{\rho^{2}}{r(r-r_{\mathrm{g}})}}\,\mathrm{d}r.
\end{equation*}

和先前从 (101.3) 到 (101.6) 所用的变换相同，此处作同样的变换，我们得到：
\begin{equation*}
  \psi_r(r)=\frac{\omega_0}{c}\int
  \sqrt{1+\frac{2r_{\mathrm{g}}}{r}-\frac{\rho^{2}}{r^{2}}}\,\mathrm{d}r.
\end{equation*}

将被积函数按 $r_{\mathrm{g}}/r$ 的幂展开，我们有：
\begin{equation*}
  \psi_r=\psi_r^{(0)}+\frac{r_{\mathrm{g}}\omega_0}{c}
  \int\frac{\mathrm{d}r}{\sqrt{r^{2}-\rho^{2}}}
  =\psi_r^{(0)}+\frac{r_{\mathrm{g}}\omega_0}{c}
  \operatorname{arch}\frac{r}{\rho},
\end{equation*}
式中 $\psi_r^{(0)}$ 相应于经典直线.

从某个非常大的距离 $R$ 到最接近中心的点 $r=\rho$ 再回到距离 $R$，光线传播过程中 $\psi_r$ 的总改变等于
\begin{equation*}
  \Delta\psi_r=\Delta\psi_r^{(0)}
  +2\frac{r_{\mathrm{g}}\omega_0}{c}\operatorname{arch}\frac{R}{\rho}.
\end{equation*}

通过对 $M=\rho\omega_0/c$ 微分得到极角 $\varphi$ 沿光线相应的改变：
\begin{equation*}
  \Delta\varphi=-\frac{\partial\Delta\psi_r}{\partial M}
  =-\frac{\partial\Delta\psi_r^{(0)}}{\partial M}
  +\frac{2r_{\mathrm{g}}R}{\rho\sqrt{R^{2}-\rho^{2}}}.
\end{equation*}

最后，过渡到极限 $R\to\infty$，并注意到直线对应于 $\Delta\varphi=\pi$，我们得到：
\begin{equation*}
  \Delta\varphi=\pi+\frac{2r_{\mathrm{g}}}{\rho}.
\end{equation*}

这意味着光线在引力场的影响下弯曲了：它的轨道是一条凹向中心的曲线（光线被“吸”向中心），其两条渐进线之间的夹角与 $\pi$ 相差
\begin{equation}
  \delta\varphi=\frac{2r_{\mathrm{g}}}{\rho}
  =\frac{4km'}{c^{2}\rho};
\end{equation}

换言之，在离场中心距离 $\rho$ 处经过的光线偏折了一个角度 $\delta\varphi$.\footnote{一条掠过太阳边缘的光线的偏折为 $\delta\varphi=1.75''$.}

\section{球形物体的引力坍缩}

在史瓦西度规 (100.14) 中，在 $r=r_{\mathrm{g}}$ 处（史瓦西球上）$g_{00}$ 变为 0，$g_{11}$ 变为无穷大.据此似可得出结论，那里必定是时空度规的奇点，因而（对给定的质量而言）不可能存在“半径”小于引力半径的物体.然而，这一结论实际上可能是不对的.从如下事实就能看出这点：行列式 $g=-r^{4}\sin^{2}\theta$ 在 $r=r_{\mathrm{g}}$ 没有任何奇异性，因而条件 $g<0$ (82.3) 并未受到破坏.下面将看到，我们正在讨论的实际上只是对于 $r<r_{\mathrm{g}}$ 建立刚性参考系的不可能性.

为了搞清时空度规在这个区域的特性，我们作如下形式的坐标变换\footnote{史瓦西奇点的物理意义是由 D. Finkelstein (1958) 用不同的变换首先解释的.度规 (102.3) 由 O. Lemaitre (1938) 首先得到.}：
\begin{equation}
  c\tau=\pm ct\pm\int\frac{f(r)\,\mathrm{d}r}{1-\dfrac{r_{\mathrm{g}}}{r}},
  \qquad
  R=ct+\int\frac{\mathrm{d}r}{\left(1-\dfrac{r_{\mathrm{g}}}{r}\right)f(r)}.
\end{equation}

于是
\begin{equation*}
  \mathrm{d}s^{2}=\frac{1-\dfrac{r_{\mathrm{g}}}{r}}{1-f^{2}}
  \left(c^{2}\mathrm{d}\tau^{2}-f^{2}\mathrm{d}R^{2}\right)
  -r^{2}(\mathrm{d}\theta^{2}+\sin^{2}\theta\,\mathrm{d}\varphi^{2}).
\end{equation*}

我们通过选择 $f(r)$ 使 $f(r_{\mathrm{g}})=1$ 来消去 $r=r_{\mathrm{g}}$ 处的奇点.如果令 $f(r)=\sqrt{r_{\mathrm{g}}/r}$，则新坐标系也是同步的（$g_{\tau\tau}=1$）.首先选择 (102.1) 中的正号，我们有：
\begin{equation*}
  R-c\tau=\int\frac{(1-f^{2})\,\mathrm{d}r}
  {\left(1-\dfrac{r_{\mathrm{g}}}{r}\right)f}
  =\int\sqrt{\frac{r}{r_{\mathrm{g}}}}\,\mathrm{d}r
  =\frac{2}{3}\frac{r^{3/2}}{r_{\mathrm{g}}^{1/2}},
\end{equation*}

或
\begin{equation}
  r=\left[\frac{3}{2}(R-c\tau)\right]^{2/3}r_{\mathrm{g}}^{1/3}
\end{equation}
（积分常数依赖于时间原点，我们令它等于零）.间隔元是：
\begin{equation}
  \mathrm{d}s^{2}=c^{2}\mathrm{d}\tau^{2}
  -\frac{\mathrm{d}R^{2}}
  {\left[\dfrac{3}{2r_{\mathrm{g}}}(R-c\tau)\right]^{2/3}}
  -\left[\frac{3}{2}(R-c\tau)\right]^{4/3}r_{\mathrm{g}}^{2/3}
  (\mathrm{d}\theta^{2}+\sin^{2}\theta\,\mathrm{d}\varphi^{2}).
\end{equation}

在这些坐标中，史瓦西球（相应于等式 $\dfrac{3}{2}(R-c\tau)=r_{\mathrm{g}}$）上的奇异性就不存在了.坐标 $R$ 处处类空，而 $\tau$ 处处类时.度规 (102.3) 是非静态的.像在每个同步参考系中一样，时间线是测地线.换言之，相对于该参考系静止的“检验”粒子是在给定场中自由运动的粒子.

对于给定的 $r$ 值，相应的世界线是 $R-c\tau=\mathrm{const}$（图 20 中的倾斜直线）.相对于参考系静止的粒子的世界线在图上显示为垂直线；粒子沿着这些线运动，在有限的固有时间隔后“落入”场中心（$r=0$），这就是度规奇点的位置.

\par\nopagebreak\noindent\begin{minipage}{\linewidth}\centering\includegraphics[width=7cm]{fig20.pdf}\\[0.3em]{\small 图 20}\end{minipage}\par\nopagebreak

我们来考虑径向光信号的传播.方程 $\mathrm{d}s^{2}=0$（对于 $\theta,\varphi=\mathrm{const}$）沿光线给出导数 $\mathrm{d}\tau/\mathrm{d}R$：
\begin{equation}
  c\frac{\mathrm{d}\tau}{\mathrm{d}R}
  =\pm\left[\frac{3}{2r_{\mathrm{g}}}(R-c\tau)\right]^{-1/3}
  =\pm\sqrt{\frac{r_{\mathrm{g}}}{r}},
\end{equation}

正负号相应于其顶点在给定世界点的光“锥”的两个边界.当 $r>r_{\mathrm{g}}$ 时（图 20 中的点 $a$）这些边界的斜率满足 $|c\,\mathrm{d}\tau/\mathrm{d}R|<1$，所以直线 $r=\mathrm{const}$（沿着它 $c\,\mathrm{d}\tau/\mathrm{d}R=1$）落入锥内.但在 $r<r_{\mathrm{g}}$ 的区域（点 $a'$）我们有 $|c\,\mathrm{d}\tau/\mathrm{d}R|>1$，所以线 $r=\mathrm{const}$，即相对于场中心静止粒子的世界线，处于锥外.锥的两个边界与线 $r=0$ 相交于有限距离处，沿着垂直线趋近于它.因为因果联系的事件不能够处于光锥以外的世界线上，所以在 $r<r_{\mathrm{g}}$ 的区域，没有粒子可以处于静止.这里所有相互作用和信号都朝着指向中心的方向传播，在有限的固有时间隔 $\tau$ 后到达那里.

类似地，在 (102.1) 中选择负号，我们会得到一个“膨胀”的参考系，其中度规与 (102.3) 的差别在于 $\tau$ 变号.它所相应的时空，在其中（$r<r_{\mathrm{g}}$ 的区域内）静止也是不可能的，但所有的信号从中心向外传播.

这里描述的结果可以应用于广义相对论中大质量物体的行为问题.

相对论性球体平衡条件的研究显示，对于一个质量足够大的物体，不能存在静平衡态（参见本教程第五卷§109）.显然，这样的物体必定会无限收缩（称为引力坍缩\footnote{这一现象的基本性质是由 J. R. Oppenheimer 和 H. Snyder 首先阐明的 (1939).}）.

在与物体没有联系、无穷远处为伽利略的参考系中（度规 (100.14)），中心物体的半径不能小于 $r_{\mathrm{g}}$.这意味着按照遥远观察者的钟，收缩物体的半径只是在 $t\to\infty$ 时渐近地趋于引力半径.容易找到这种依赖关系的极限形式.

收缩物体表面上的粒子在所有时间内都处于恒定质量 $m$（物体的总质量）的引力场中.当 $r\to r_{\mathrm{g}}$ 时引力变得非常大；但物体的密度（随之压强）仍然是有限的，因此我们可以忽略压力，而把确定物体半径同时间的依赖关系 $r=r(t)$，化为考虑质量 $m$ 的场中检验粒子的自由下落.

对于史瓦西场中的下落而言，函数 $r(t)$ 由积分 (101.4) 给出，由于运动是纯径向的，那里的角动量 $M=0$.因此，若下落在某个时刻 $t_0$ 以零速度在离中心的“距离” $r_0$ 处开始，粒子的能量是 $\mathcal{E}_0=mc^{2}\sqrt{1-r_{\mathrm{g}}/r_0}$，而在时刻 $t$ 达到“距离” $r$，我们有：
\begin{equation}
  c(t-t_0)=\sqrt{1-\frac{r_{\mathrm{g}}}{r_0}}\int_{r}^{r_0}
  \frac{\mathrm{d}r}{\left(1-\dfrac{r_{\mathrm{g}}}{r}\right)
  \sqrt{\dfrac{r_{\mathrm{g}}}{r}-\dfrac{r_{\mathrm{g}}}{r_0}}}.
\end{equation}

这个积分在 $r\to r_{\mathrm{g}}$ 时像 $-r_{\mathrm{g}}\ln(r-r_{\mathrm{g}})$ 一样发散.因此我们得到趋于 $r_{\mathrm{g}}$ 的渐进公式：
\begin{equation}
  r-r_{\mathrm{g}}=\mathrm{const}\cdot\mathrm{e}^{-ct/r_{\mathrm{g}}}.
\end{equation}

因此，物体在趋向引力半径的最后阶段按照指数律坍缩，其特征时间非常小（$\sim r_{\mathrm{g}}/c$）.

尽管从外面观察收缩速率渐进地趋向于零，但粒子以其固有时测量的下落速度 $v$ 会增加并趋于光速.实际上，按照定义 (88.10)：
\begin{equation*}
  v^{2}=\left(\frac{\sqrt{-g_{11}}\,\mathrm{d}r}
  {\sqrt{g_{00}}\,\mathrm{d}t}\right)^{2}.
\end{equation*}

从 (100.14) 取来 $g_{11}$ 和 $g_{00}$，从 (102.5) 取来 $\mathrm{d}r/\mathrm{d}t$，我们得到：
\begin{equation}
  1-\frac{v^{2}}{c^{2}}
  =\frac{1-\dfrac{r_{\mathrm{g}}}{r}}{1-\dfrac{r_{\mathrm{g}}}{r_0}}.
\end{equation}

按照外面观察者的钟，趋向引力半径要花无限长的时间，但它却只占用有限的固有时（即随物体一起运动的参考系中的时间）间隔.这一点从上面所作的一般分析中已经很清楚了，但我们也可以通过计算作为不变积分
\begin{equation*}
  c\tau=\int\mathrm{d}s
  =\int\left[c^{2}g_{00}\left(\frac{\mathrm{d}t}{\mathrm{d}r}\right)^{2}
  +g_{11}\right]^{1/2}\mathrm{d}r
\end{equation*}
的固有时 $\tau$ 来直接予以验证.从 (102.5) 取下落粒子的 $\mathrm{d}r/\mathrm{d}t$，我们得到从 $r_0$ 下落到 $r$ 的固有时：
\begin{equation}
  \tau-\tau_0=\frac{1}{c}\int_{r}^{r_0}
  \left(\frac{r_{\mathrm{g}}}{r}-\frac{r_{\mathrm{g}}}{r_0}\right)^{-1/2}
  \mathrm{d}r.
\end{equation}

这个积分在 $r\to r_{\mathrm{g}}$ 时收敛.

（按固有时测量）达到引力半径后，物体将继续收缩，带着它所有的粒子在有限时间内抵达中心；每部分物质坍缩进中心的时刻是时空度规的真正奇点.不过，我们根本观察不到史瓦西球内部物体坍缩的这一过程.物体穿越该球表面的时间相应于 $t=\infty$；我们可以说，对于遥远的观察者，坍缩进史瓦西球的整个过程发生在“无限长的时间以后”——这是时间相对性的一个极端例子.这幅图像中当然没有任何逻辑矛盾.与此完全相应的是上面关于收缩坐标系性质的论断：在该系中没有信号从史瓦西球内出来.粒子或光线只能沿着一个方向——向内——同这个球相交（在共动参考系中），一旦穿过那里，就不能再出来了.这个“单向阀”称为事件视界.

对于外面的观察者来说，收缩到引力半径是通过物体的“自关闭”实现的.从物体送出信号的传播时间趋于无穷大：对于光信号 $c\,\mathrm{d}t=\mathrm{d}r/(1-r_{\mathrm{g}}/r)$，从 $r$ 传到某个 $r_0>r$ 的时间由如下积分给出：
\begin{equation}
  c\Delta t=\int_{r}^{r_0}\frac{\mathrm{d}r}{1-\dfrac{r_{\mathrm{g}}}{r}}
  =r_0-r+r_{\mathrm{g}}\ln\frac{r_0-r_{\mathrm{g}}}{r-r_{\mathrm{g}}},
\end{equation}

它像积分 (102.5) 一样当 $r\to r_{\mathrm{g}}$ 时发散.同无限远观察者的时间 $t$ 的间隔相比，物体表面固有时的间隔按如下比例缩短了：
\begin{equation*}
  \sqrt{g_{00}}=\sqrt{1-\frac{r_{\mathrm{g}}}{r}};
\end{equation*}

因而当 $r\to r_{\mathrm{g}}$ 时，物体上所有的过程对于外面的观察者显得被“冻结”.由物体发出而被遥远观察者收到的谱线的频率减小，但这不仅是引力红移的效应，而且也是由于源运动的多普勒频移效应，这个源正同球面一起落向中心.当球的半径已接近 $r_{\mathrm{g}}$（因此下落速度已接近光速）时，这个效应将频率减小如下因子：
\begin{equation*}
  \frac{\sqrt{1-v^{2}/c^{2}}}{1+v/c}
  \approx\frac{1}{2}\sqrt{1-\frac{v^{2}}{c^{2}}}.
\end{equation*}

在这两个效应的影响下，观测频率随 $r\to r_{\mathrm{g}}$ 按如下规律趋于零：
\begin{equation}
  \omega=\mathrm{const}\left(1-\frac{r_{\mathrm{g}}}{r}\right).
\end{equation}

因此，从遥远观察者看来，引力坍缩导致“被冻结”物体的外貌，它不向周围空间发出任何信号，而只通过其静态引力场同外部世界作用.这样的结构称为黑洞或坍缩星.

最后我们再作一个方法论性质的评论.我们已经看到，对于真空中的中心场，“外部观察者的参考系”（它在无穷远处是惯性系）并不完备：其中没有在史瓦西球内运动粒子的世界线的容身之处.度规 (102.3) 在史瓦西球内仍然适用，但这个参考系在某种意义上也不完备.在这个系统中，考虑一个粒子从中心向外作径向运动.当 $\tau\to\infty$ 时，它的世界线向外走到无穷远，而在 $\tau\to-\infty$ 时，必须渐近地趋于 $r=r_{\mathrm{g}}$，因为在这个度规中，在史瓦西球内部，运动只能沿朝着中心的方向进行.另一方面，粒子从 $r=r_{\mathrm{g}}$ 出射到任何 $r>r_{\mathrm{g}}$ 的给定点发生在有限的固有时间隔内.因而按固有时，粒子必须先在内部趋向史瓦西球，然后才能开始在其外部运动；但粒子历史的这一部分并没有被这个特别的参考系包含.\footnote{构造没有这种不完备性的参考系将在下节末考虑.}

我们强调指出，这种不完备性只是产生于对场的度规的形式处理，那里的场被看做是由质点产生的.在实际物理问题中，例如有广延物体的坍缩，这种不完备性并不存在：将度规 (102.3) 同物质内部的解吻合而得到的解当然会是完备的，它将描述粒子所有可能运动的整个历史（在 $r>r_{\mathrm{g}}$ 区域朝向中心运动粒子的世界线必须从球面开始，甚至在它收缩到史瓦西球内之前）.

\begin{xiti}

\noindent{\heiti 习题1}\quad 对于黑洞的场中的粒子，求其圆轨道的半径 (S. A. Kaplan, 1949).

解：对于在史瓦西场中运动的粒子，依赖关系 $r=r(t)$ 由 (101.4) 给出，或者以微分形式：
\begin{equation}
  \frac{1}{1-r_{\mathrm{g}}/r}\frac{\mathrm{d}r}{c\,\mathrm{d}t}
  =\frac{1}{\mathcal{E}_0}
  \left[\mathcal{E}_0^{2}-U^{2}(r)\right]^{1/2},
  \tag{1}
\end{equation}
式中
\begin{equation*}
  U(r)=mc^{2}\left[\left(1-\frac{r_{\mathrm{g}}}{r}\right)
  \left(1+\frac{M^{2}}{m^{2}c^{2}r^{2}}\right)\right]^{1/2}
\end{equation*}
（$m$ 是粒子的质量，$r_{\mathrm{g}}=2km'/c^{2}$ 是质量为 $m'$ 的中心天体的引力半径）.$U(r)$ 在下述意义上起着“有效势能”的作用：条件 $\mathcal{E}_0\geqslant U(r)$ 决定运动的容许范围（类似于非相对性理论）.图 21 显示了 $U(r)$ 对于粒子角动量 $M$ 的各种值的曲线.

\par\nopagebreak\noindent\begin{minipage}{\linewidth}\centering\includegraphics[width=8cm]{fig21.pdf}\\[0.3em]{\small 图 21}\end{minipage}\par\nopagebreak

轨道的曲率半径及 $\mathcal{E}_0$ 和 $M$ 的相应值由函数 $U(r)$ 的极值决定，极小值相应于稳定轨道，极大值相应于不稳定轨道.同时解方程 $U(r)=\mathcal{E}_0,\ U'(r)=0$ 给出：
\begin{equation*}
  \frac{r}{r_{\mathrm{g}}}
  =\frac{M^{2}}{m^{2}c^{2}r_{\mathrm{g}}^{2}}
  \left[1\pm\sqrt{1-\frac{3m^{2}c^{2}r_{\mathrm{g}}^{2}}{M^{2}}}\right],
  \qquad
  \mathcal{E}_0=Mc\sqrt{\frac{2}{rr_{\mathrm{g}}}
  \left(1-\frac{r_{\mathrm{g}}}{r}\right)},
\end{equation*}
式中根号前的正号指稳定轨道，负号指不稳定轨道.最靠近中心的稳定轨道具有参数
\begin{equation*}
  r=3r_{\mathrm{g}},\qquad
  M=\sqrt{3}\,mcr_{\mathrm{g}},\qquad
  \mathcal{E}_0=\sqrt{8/9}\,mc^{2}.
\end{equation*}

不稳定轨道的最小半径是 $3r_{\mathrm{g}}/2$，在极限 $M\to\infty$，$\mathcal{E}_0\to\infty$ 下达到.图 22 描绘了 $r/r_{\mathrm{g}}$ 对于 $M/(mcr_{\mathrm{g}})$ 的依赖关系；上半支给出稳定轨道的半径，下半支给出不稳定轨道的半径.\footnote{为了比较我们回想起，在牛顿场中圆轨道在离中心任何距离处都是可能（而且稳定）的，半径与角动量由公式 $r=M^{2}/(km'm^{2})$ 联系起来.}

\par\nopagebreak\noindent\begin{minipage}{\linewidth}\centering\includegraphics[width=7cm]{fig22.pdf}\\[0.3em]{\small 图 22}\end{minipage}\par\nopagebreak

\noindent{\heiti 习题2}\quad 对于同一场中的运动，求来自无穷远的 a) 非相对论的，b) 极端相对论的，粒子引力俘获的截面 (Ya. B. Zel'dovich 和 I. D. Novikov, 1964).

解：(a) 对于非相对论速度 $v_{\infty}$（无穷远处），粒子的能量是 $\mathcal{E}_0\approx mc^{2}$.从图 21 可见，直线 $\mathcal{E}_0=mc^{2}$ 处于角动量 $M<2mcr_{\mathrm{g}}$，即碰撞参量 $\rho<2cr_{\mathrm{g}}/v_{\infty}$ 的所有势能曲线的上方.具有这样 $\rho$ 值的所有粒子将遭到引力俘获：它们（当 $t\to\infty$ 时渐近地）达到史瓦西球，而不再出射到无穷远.俘获截面是
\begin{equation*}
  \sigma=4\pi r_{\mathrm{g}}^{2}\left(\frac{c}{v_{\infty}}\right)^{2}.
\end{equation*}

(b) 通过代换 $m\to0$，将习题1 的方程 (1) 过渡到极端相对论粒子（或光线）情形.再引入碰撞参量 $\rho=cM/\mathcal{E}_0$，我们得到：
\begin{equation*}
  \frac{1}{1-r_{\mathrm{g}}/r}\frac{\mathrm{d}r}{c\,\mathrm{d}t}
  =\sqrt{1-\frac{\rho^{2}}{r^{2}}
  +\frac{\rho^{2}r_{\mathrm{g}}}{r^{3}}}.
\end{equation*}

径向运动的边界（转折点）由根号内表达式的根决定.图 23 绘出了它们作为 $\rho$ 的函数；可能的运动范围相应于平面上没有阴影的部分.该曲线的极小值在点
\begin{equation*}
  \rho=\frac{3\sqrt{3}}{2}r_{\mathrm{g}},\qquad
  r=\frac{3}{2}r_{\mathrm{g}}.
\end{equation*}

对于较小的碰撞参量值，粒子到不了转折点，即它将朝史瓦西球运动.于是我们得到俘获截面
\begin{equation*}
  \sigma=\frac{27}{4}\pi r_{\mathrm{g}}^{2}.
\end{equation*}

\par\nopagebreak\noindent\begin{minipage}{\linewidth}\centering\includegraphics[width=6.5cm]{fig23.pdf}\\[0.3em]{\small 图 23}\end{minipage}\par\nopagebreak

\end{xiti}

\section{类尘埃球的坍缩}

要讨论坍缩物体内部状态改变的过程（包括处理它在史瓦西球内压缩过程），需要求解物质介质内部引力场的爱因斯坦方程.在中心对称情形下，如果我们忽略物质的压强：$p=0$ (R. Tolman, 1934)，可以求得场方程一般形式的解.尽管在真实情况下作近似常常是容许的，但这个问题的通解还是有相当大的方法论意义.

正如§97中所示，对于类尘埃介质，我们可以选择既是同步也是共动的参考系.\footnote{物质必须“无旋转”地运动（参见§97第三个脚注），在本例中，这一条件一定满足，因为球对称意味着物质的纯径向运动.}用 $\tau$ 和 $R$ 表示这样选择的时间和径向坐标，我们写出如下形式的球对称线元：\footnote{在本节中，令 $c=1$.}
\begin{equation}
  \mathrm{d}s^{2}=\mathrm{d}\tau^{2}
  -\mathrm{e}^{\lambda(\tau,R)}\mathrm{d}R^{2}
  -r^{2}(\tau,R)(\mathrm{d}\theta^{2}
  +\sin^{2}\theta\,\mathrm{d}\varphi^{2}).
\end{equation}

函数 $r(\tau,R)$ 是如此定义的“半径”，使 $2\pi r$ 是圆（中心在原点）的周长.形式 (103.1) 唯一地确定了 $\tau$ 的选择，但仍然容许形如 $R=R(R')$ 的任意径向坐标变换.

这个度规的里奇张量分量的计算导致如下爱因斯坦方程组\footnote{比较§100习题5. 如果令 $\nu=0,\ \mathrm{e}^{\mu}=r^{2},\ p=0$，从该习题的方程 (2)—(5) 可分别得到方程 (103.2)—(103.5). 我们注意到，当 $p=0$ 时，同题方程 (6) 的第二个方程给出 $\nu'=0$，即 $\nu=\nu(\tau)$；因而在选择 $\tau$ 时，度规 (1) 中留下的任意性容许我们令 $\nu=0$，这就再次证明了引入同步共动参考系的可能性.}：
\begin{equation}
  -\mathrm{e}^{-\lambda}r'^{2}+2r\ddot{r}+\dot{r}^{2}+1=0,
\end{equation}

\begin{equation}
  -\frac{\mathrm{e}^{-\lambda}}{r}(2r''-r'\lambda')
  +\frac{\dot{r}\dot{\lambda}}{r}
  +\ddot{\lambda}+\frac{\dot{\lambda}^{2}}{2}
  +\frac{2\ddot{r}}{r}=0,
\end{equation}

\begin{equation}
  -\frac{\mathrm{e}^{-\lambda}}{r^{2}}
  (2rr''+r'^{2}-rr'\lambda')
  +\frac{1}{r^{2}}(r\dot{r}\dot{\lambda}+\dot{r}^{2}+1)
  =8\pi k\varepsilon,
\end{equation}

\begin{equation}
  2\dot{r}'-\dot{\lambda}r'=0,
\end{equation}
这里撇表示对 $R$ 微分，点表示对 $\tau$ 微分.

方程 (103.5) 直接对时间积分，得出
\begin{equation}
  \mathrm{e}^{\lambda}=\frac{r'^{2}}{1+f(R)},
\end{equation}
式中 $f(R)$ 是任意函数，只遵从条件 $1+f>0$.将这个表达式代入 (103.2)，我们得到
\begin{equation*}
  2r\ddot{r}+\dot{r}^{2}-f=0
\end{equation*}
（代入 (103.3) 得不到任何新结果）.这个方程的初积分是
\begin{equation}
  \dot{r}^{2}=f(R)+\frac{F(R)}{r},
\end{equation}
式中 $F(R)$ 是另一个任意函数.因此
\begin{equation*}
  \tau=\pm\int\frac{\mathrm{d}r}
  {\sqrt{f+\dfrac{F}{r}}}.
\end{equation*}

从该积分得到的函数 $r(\tau,R)$ 可以写成参数形式：
\begin{equation}
  r=\frac{F}{2f}(\cosh\eta-1),\qquad
  \tau_0(R)-\tau=\frac{F}{2f^{3/2}}(\sinh\eta-\eta),
  \quad\text{当}\ f>0,
\end{equation}

\begin{equation}
  r=\frac{F}{-2f}(1-\cos\eta),\qquad
  \tau_0(R)-\tau=\frac{F}{2(-f)^{3/2}}(\eta-\sin\eta),
  \quad\text{当}\ f<0,
\end{equation}

式中 $\tau_0(R)$ 还是一个任意函数.如果 $f=0$，则
\begin{equation}
  r=\left(\frac{9F}{4}\right)^{1/3}
  \left[\tau_0(R)-\tau\right]^{2/3},
  \quad\text{当}\ f=0.
\end{equation}

在所有情形下，将 (103.6) 代入 (103.4) 并用 (103.7) 消去 $f$，我们得到物质密度的如下表达式\footnote{函数 $F,f,\tau_{0}$ 只满足假设 $\mathrm{e}^{\lambda},r$ 和 $\varepsilon$ 的正性条件. 加上上面给出的条件 $1+f>0$，推得 $f>0$. 我们也假设 $F'>0,\ r'>0$；这就排除了导致物质在其径向运动中穿过球层的情况.}：
\begin{equation}
  8\pi k\varepsilon=\frac{F'}{r'r^{2}}.
\end{equation}

公式 (103.6)—(103.11) 确定了要求的通解\footnote{然而这并不包括如下特例，即 $r=r(\tau)$，不依赖于 $R$，故 (103.5) 化为恒等式；参见 V. A. Ruban, JETP, 56, 1914 (1969)；Soviet Phys. JETP, 29, 1027 (1969). 不过，这种情况并不对应有限物体坍缩问题的条件.}.我们发现，它只依赖于两个“物理上不同”的任意函数：尽管其中出现了三个函数 $f,F,\tau_0$，坐标 $R$ 还可经受任意变换 $R=R(R')$.这个数字正好对应于：物质最一般的中心对称分布由两个函数（物质的密度分布和径向速度）给出，而中心对称的自由引力场并不存在.

因为参考系随物质运动，每个物质粒子对应于确定的 $R$ 值；对于这个 $R$ 值，函数 $r(\tau,R)$ 决定了给定粒子的运动规律，而导数 $\dot{r}$ 是其径向速度.这里得到的解的重要性质是，在从 0 到某个 $R_0$ 的区间上给定出现于解中的任意函数，就完全决定了这个半径的球的行为；它并不依赖于在 $R>R_0$ 时如何给定这些函数.我们可以自动获得对于任何有限球内部问题的解.球的总质量，按 (100.23)，由如下积分给定：
\begin{equation*}
  m=4\pi\int_{0}^{r(\tau,R_0)}\varepsilon r^{2}\,\mathrm{d}r
  =4\pi\int_{0}^{R_0}\varepsilon r^{2}r'\,\mathrm{d}R.
\end{equation*}

代入 (103.11) 并注意到 $F(0)=0$（当 $R=0$，必有 $r=0$），我们得到
\begin{equation}
  m=\frac{F(R_0)}{2k},\qquad r_{\mathrm{g}}=F(R_0)
\end{equation}
（$r_{\mathrm{g}}$ 是球的引力半径）.

对于 $F=\mathrm{const}\neq0$，我们从 (103.11) 得 $\varepsilon=0$，所以这个解适用于真空，即它描述质点的场（位于中心，度规的奇点）.所以，令 $F=r_{\mathrm{g}},\ f=0,\ \tau_0=R$，我们得到度规 (102.3)\footnote{$F=0$ 的情况（那里 (103.7) 给出 $r=\sqrt{f}\,(\tau-\tau_{0})$）对应于场不存在；通过简单的变量代换，该度规可以化为伽利略形式.}.

公式 (103.8)—(103.10) 描述了球的收缩和膨胀（依赖于参量 $\eta$ 的取值范围）；对于场方程来说，两者都是同样容许的.不稳定的大质量物体行为的重要问题对应于收缩——引力坍缩.解 (103.8)—(103.10) 描绘了这样的图景：当 $\tau$ 增加到趋于 $\tau_0$ 时发生收缩.时刻 $\tau=\tau_0(R)$ 对应于到达具有给定径向坐标 $R$ 的物质的中心（在那里我们必须有 $\tau_0'>0$）.

当 $\tau\to\tau_0(R)$ 时球内度规的极限特征对所有三种情形 (103.8)—(103.10) 都相同：
\begin{equation}
  r\approx\left(\frac{9F}{4}\right)^{1/3}(\tau_0-\tau)^{2/3},
  \qquad
  \mathrm{e}^{\lambda/2}\approx\left(\frac{2F}{3}\right)^{1/3}
  \frac{\tau_0'}{\sqrt{1+f}}(\tau_0-\tau)^{-1/3}.
\end{equation}

这意味着，所有径向距离（在这里所用的共动参考系内）趋向无穷大，而切向距离则趋于零（像 $\tau-\tau_0$ 那样）\footnote{经过中心的“平面”上的几何是，可能存在一个旋转锥面，随时间沿其母线伸长，同时沿其所有的圆周收缩.}.物质密度相应地无限增加\footnote{在这个解中任何质量的球都发生引力坍缩是忽略压强的自然结果. 显然，当 $\varepsilon\to\infty$ 时，从物理观点看，假设物质类尘埃是绝不容许的，我们应当用极端相对论的物态方程 $p=\varepsilon/3$. 然而，看起来极限压缩规律的一般特征在很大程度上与物态方程无关，参见 E. M. Lifshitz 和 I. M. Khalatnikov, JETP 39, 149, 1960；Soviet Phys. JETP, 12, 108, 1961.}：
\begin{equation}
  8\pi k\varepsilon\approx
  \frac{2F'}{3F\tau_0'(\tau_0-\tau)}.
\end{equation}

因此，与我们在§102中的评论一致，整个物质分布都坍缩到中心\footnote{$\tau_{0}=\mathrm{const}$ 的情况，包含了一个特例，即完全均匀球的坍缩，参见习题.}.

在函数 $\tau_0(R)=\mathrm{const}$（即所有粒子同时到达中心）的特殊情形下，收缩球内部的度规具有不同的特征.在这种情形下，
\begin{equation}
\begin{aligned}
  &r\approx\left(\frac{9F}{4}\right)^{1/3}(\tau_0-\tau)^{2/3},
  \qquad
  \mathrm{e}^{\lambda/2}\approx\left(\frac{2}{3}\right)^{1/3}
  \frac{F'}{2F^{2/3}\sqrt{1+f}}(\tau_0-\tau)^{2/3},\\
  &8\pi k\varepsilon\approx\frac{4}{3(\tau_0-\tau)^{2}}.
\end{aligned}
\end{equation}

即，随着 $\tau\to\tau_0$，所有距离（无论径向和切向）都按同样规律 $\sim(\tau_0-\tau)^{2/3}$ 趋于零；物质密度像 $(\tau_0-\tau)^{-2}$ 那样趋于无穷大，而且在这个极限，物质的分布变得均匀.

我们注意到，在所有情形，坍缩球表面穿过史瓦西球（$r(\tau,R_0)=r_{\mathrm{g}}$）的时刻，对于其内部的动力学（由共动参考系的度规描述）并不重要.然而，在每一时刻，该球的一定部分已经在自己的“事件视界”以内.正如 $F(R_0)$，通过 (103.12)，决定了整个球的引力半径一样，对于任何给定的 $R$ 值，$F(R)$ 是该球在 $R=\mathrm{const}$ 的球面内那一部分的引力半径；因此，球的这一部分在每一时刻 $\tau$ 由条件 $r(\tau,R)\leqslant F(R)$ 决定.

最后，我们来说明如何能够用这些公式求解§102末提出的问题：为质点的场构造最完备的参考系\footnote{这样的参考系首先是克鲁斯卡尔（M. Kruskal, 1960）用其他变量找到的（参见 Phys. Rev. 119, 1743, 1960）. 这里给出的解的形式出自 I. D. Novikov, 1963，用的是同步参考系.}.

为了实现这一目的，必须从能够包含收缩和膨胀两个时空区域的真空度规出发.方程 (103.8) 就是这样的解，其中我们必须令 $F=\mathrm{const}=r_{\mathrm{g}}$.再选择
\begin{equation*}
  f=-\frac{1}{(R/r_{\mathrm{g}})^{2}+1},\qquad
  \tau_0=\frac{\pi}{2}r_{\mathrm{g}}(-f)^{-3/2},
\end{equation*}
我们得到：
\begin{equation}
\begin{aligned}
  &\frac{r}{r_{\mathrm{g}}}
  =\frac{1}{2}\left(\frac{R^{2}}{r_{\mathrm{g}}^{2}}+1\right)
  (1-\cos\eta),\\
  &\frac{\tau}{r_{\mathrm{g}}}
  =\frac{1}{2}\left(\frac{R^{2}}{r_{\mathrm{g}}^{2}}+1\right)^{3/2}
  (\pi-\eta+\sin\eta);
\end{aligned}
\end{equation}
式中参量 $\eta$ 取值从 0 到 $2\pi$，时间 $\tau$（对于给定的 $R$）单调减小，而 $r$ 从零增加，经过一个极大值，然后再降到零.

在图 24 中，曲线 $ACB$ 和 $A'C'B'$ 相应于点 $r=0$（参数值 $\eta=2\pi$ 和 $\eta=0$）.曲线 $AOA'$ 和 $BOB'$ 相应于史瓦西球 $r=r_{\mathrm{g}}$.$A'C'B'$ 和 $A'OB'$ 之间的时空区域只可能有从中心向外的运动，而 $ACB$ 和 $AOB$ 之间是只发生向中心运动的区域.

\par\nopagebreak\noindent\begin{minipage}{\linewidth}\centering\includegraphics[width=8cm]{fig24.pdf}\\[0.3em]{\small 图 24}\end{minipage}\par\nopagebreak

相对于这个参考系静止的粒子的世界线是垂线（$R=\mathrm{const}$）.它从 $r=0$（点 $a$）出发，在点 $b$ 同史瓦西球相交，在时刻 $\tau=0$ 达到其最远距离 $\left[r=r_{\mathrm{g}}\left(\dfrac{R^{2}}{r_{\mathrm{g}}^{2}}+1\right)\right]$，此后粒子开始向史瓦西球下落，在点 $c$ 穿过它，再次达到 $r=0$（点 $d$）的时间是
\begin{equation*}
  \tau=r_{\mathrm{g}}\frac{\pi}{2}
  \left(\frac{R^{2}}{r_{\mathrm{g}}^{2}}+1\right)^{3/2}.
\end{equation*}

这个参考系是完备的：在场中运动的任何粒子的世界线的两端要么在真奇点 $r=0$，要么在无穷远.不完备的度规 (102.3) 只覆盖到 $AOA'$ 右边（或 $BOB'$ 左边）的区域，而“膨胀”参考系覆盖到 $BOB'$ 右边（或 $AOA'$ 左边）的区域.有度规 (100.14) 的史瓦西参考系只覆盖到 $BOA'$ 右边（或 $AOB'$ 左边）的区域.

\begin{xiti}

\noindent{\heiti 习题}\quad 求类尘埃均匀球（其物质开始时处于静止）引力坍缩问题的内部解.

解：令
\begin{equation*}
  \tau_0=\mathrm{const},\qquad
  f=-\sin^{2}R,\qquad
  F=2a_0\sin^{3}R,
\end{equation*}
我们得到
\begin{equation}
  r=a_0\sin R(1-\cos\eta),\qquad
  \tau-\tau_0=a_0(\eta-\sin\eta)
  \tag{1}
\end{equation}
（径向坐标 $R$ 是无量纲的，取值从 0 到 $2\pi$），则密度是
\begin{equation}
  8\pi k\varepsilon=\frac{6}{a_0^{2}(1-\cos\eta)^{3}}.
  \tag{2}
\end{equation}

对于给定的 $\tau$，它与 $R$ 无关，所以球是均匀的.由 (1) 我们把带 $r$ 的度规 (103.1) 表为形式
\begin{equation}
\begin{aligned}
  &\mathrm{d}s^{2}=\mathrm{d}\tau^{2}
  -a^{2}(\tau)\left[\mathrm{d}R^{2}
  +\sin^{2}R(\mathrm{d}\theta^{2}
  +\sin^{2}\theta\,\mathrm{d}\varphi^{2})\right],\\
  &a=a_0(1-\cos\eta).
\end{aligned}
\tag{3}
\end{equation}
我们注意到，它与完全充满均匀类尘埃物质的宇宙的度规，即弗里德曼解一致（§112）——这是一个完全自然的结果，因为从均匀分布物质中切割出来的球具有中心对称性\footnote{度规 (3) 相应于常正曲率空间. 以类似的方式，若令 $f=\sinh^{2}R,\ F=2a_{0}\sinh^{3}R$，我们得到相应于常负曲率空间的解（§113）.}.

对于选定的常数 $a_0,\tau_0$，解 (1) 满足开始提出的条件.为方便起见，对参量作改变（$\eta\to\pi-\eta$），我们将解写为形式
\begin{equation}
  r=\frac{r_0}{2}\frac{\sin R}{\sin R_0}(1+\cos\eta),
  \qquad
  \tau=\frac{r_0}{2\sin R_0}(\eta+\sin\eta),
  \tag{4}
\end{equation}
这里（按照 (103.12)）球的引力半径是 $r_{\mathrm{g}}=r_0\sin^{2}R_0$.在初始时刻（$\tau=0$，$\eta=0$）物质是静止的（$\dot{r}=0$），而 $2\pi r_0=2\pi r(0,R_0)$ 是球的初始周长.该物质坍缩入中心的时刻是 $\tau=\pi r_0/(2\sin R_0)$.

遥远观察者参考系（史瓦西系）中的时间 $t$ 和球上的固有时 $\tau$ 的关系由如下方程表示：
\begin{equation*}
  \mathrm{d}\tau^{2}=\left(1-\frac{r_{\mathrm{g}}}{r}\right)\mathrm{d}t^{2}
  -\frac{\mathrm{d}r^{2}}{1-r_{\mathrm{g}}/r},
\end{equation*}
式中 $r$ 指相应于球表面的值 $r(\tau,R_0)$.由该式的积分得到 $t$ 作为同一参量 $\eta$ 的函数的如下表达式：
\begin{equation}
  \frac{t}{r_{\mathrm{g}}}
  =\ln\frac{\cot R_0+\tan\dfrac{\eta}{2}}
  {\cot R_0-\tan\dfrac{\eta}{2}}
  +\cot R_0\left[\eta
  +\frac{1}{2\sin^{2}R_0}(\eta+\sin\eta)\right]
  \tag{5}
\end{equation}
（这里时刻 $t=0$ 相应于 $\tau=0$）.球的表面穿过史瓦西球（$r(\tau,R_0)=r_{\mathrm{g}}$）对应参量 $\eta$ 的值由如下方程确定：
\begin{equation*}
  \cos^{2}\frac{\eta}{2}=\frac{r_{\mathrm{g}}}{r_0}
  =\sin^{2}R_0.
\end{equation*}

当趋近这个值时，时间 $t$ 趋于无穷大，与§102中的论断一致\footnote{由 (4) 式决定的函数 $r(\tau,R_{0})$ 当然与由外部度规计算并由积分 (102.8) 给出的函数一致. 同样的论断适用于由 (4) 和 (5) 决定的函数 $t(r)$；它与积分 (102.5) 给出的函数一致.}.

\end{xiti}
\section{非球形转动物体的引力坍缩}

严格来说，前面两节的所有结果只适用于严格球对称的物体.不过，简单的论证显示，对于那些稍许偏离球对称的物体，引力坍缩的定性图景仍然相同 (A. G. Doroshkevich, Ya. B. Zel'dovich 和 I. D. Novikov, 1965).

首先我们来考虑这样的物体，它们偏离中心对称性与物质的分布有关，而与物体的整体转动无关.

很明显，既然大质量中心对称的物体是引力不稳定的，这种不稳定对于对称性的小扰动仍将保留，所以这样的物体将会坍缩.将弱不对称性处理成小扰动，我们可以追踪它在物体收缩过程中（在共动系内）的发展.一般说来，扰动随物体密度的增加而增加.但如果扰动在收缩开始时足够弱，它们在物体达到其引力半径时仍将很小；在§103曾经指出，这个时刻对于收缩物体内部的动力学毫无特殊之处，而它的密度当然仍是有限的\footnote{在§115里将处理不稳定、无边界的均匀物质分布中扰动的发展（那里得到的公式可以同样应用于膨胀和收缩两种情况）. 未扰动或无界延伸物体的非均匀性并不改变这里做出的结论.}.

由于物体内部的扰动很小，它在外部产生的中心对称引力场的扰动依然也很小.这意味着“事件视界”，即史瓦西球的表面也几乎保持不变，没有任何东西能够阻止坍缩物体穿过它（在共动参考系内）.

因为没有信号从事件视界以内送出，所以没有关于物体内部扰动进一步发展的信息到达外部观察者；整个过程对于遥远观察者仍然是“无限延迟”的.由此可知，当物体渐近地趋于引力半径时，坍缩物体的引力场相对于外部参考系趋向于变成稳态.这个过程的特征时间非常短（$\sim r_{\mathrm{g}}/c$），在这段期间我们可以假设，外部空间只存在那些在中心对称引力场中较早发展起来的扰动.但是随着时间的进程，所有的扰动都必定作为引力波耗散进空间，走向无限远（或者穿过视界）.

在形成中的黑洞的外部引力场里，也不能保留与时间无关的静态扰动.这个结论可以从应用于真空史瓦西球中恒定扰动的分析得到.这样的分析显示，在静态情形下，每一（在无穷远处减弱的）扰动，在趋近未扰动问题的史瓦西球面时会无限制地增加\footnote{参见 T. Regge 和 J. A. Wheeler, Phys. Rev. 108, 1063 (1957). 这里要强调，我们说的是来自中心体本身的扰动. 施予无穷远的条件排除了静态扰动来自外源的情况：在这样的情况下，小扰动只给史瓦西球一些扭曲，并不改变它的定性性质，也不在其上产生真正的时空奇点.}；但是，正如已经说过的那样，在目前情况中没有外部场出现大扰动的理由.

物体密度分布中对于球对称的偏离，可以用该分布的四极矩和更高的多极矩描述；其中每一个都对外部引力场作出了自己的贡献.我们的论断意味着，外场中所有这样的扰动在坍缩的最后阶段（从外面观察者的观点看）都会渐次减弱\footnote{关于这种衰减律，参见 R. H. Price, Phys. Rev. D, 5, 2419, 2439 (1972). 坍缩期间，外引力场初始静态的 $l$ 极扰动的衰减律是 $1/t^{2l+2}$.}.产生的黑洞引力场依然是中心对称的史瓦西场，只由黑洞的总质量决定.

物体坍缩到其事件视界以内（从外面的参考系不可观测）的终极命运问题还不完全清楚.看来我们可以断言，这里坍缩也会终止于时空度规的真奇点，但奇点的类型与中心对称的情况完全不同.然而，这个问题目前尚未完全搞清楚.

让我们回到球对称弱扰动不是与密度分布，而是与物体整体转动相关的情况；弱扰动的假设意味着转动必须足够慢.除了一个例外，所有以前的论断仍然有效.从一开始就很清楚，由于物体的总角动量 $M$ 守恒，黑洞的场不能只依赖于物体的质量，与此相应的是，在中心对称引力场的不含时稳态（而非静态）扰动中，有一个扰动不在 $r\to r_{\mathrm{g}}$ 时无限增加.这正好是与物体的转动相关的扰动，可通过在史瓦西度规张量 $g_{ik}$（在坐标 $x^0=t,\ x^1=r,\ x^2=\theta,\ x^3=\varphi$ 中）加上很小的非对角分量\footnote{在本节中令 $c=1$.}
\begin{equation}
  g_{03}=\frac{2kM}{c^{2}r}\sin^{2}\theta
\end{equation}
予以描述（参见§105习题）.当物体趋于引力半径时，这个表达式仍然有效（在外部空间中），因此缓慢转动黑洞的引力场（在小角动量 $M$ 的情形准确到一级近似），将是中心对称的史瓦西场加上小修正 (104.1).这个场不再是静态的而只是稳态的.

如果引力坍缩对于球对称的小扰动是可能的，那么如果在某个有限范围之内偏离球形，必定也可能发生同样性质的坍缩（连同物体运动到事件视界以内）；决定这一范围的条件尚未建立.无论这些条件如何，看来我们可以断言，从外部观察者的观点来看，由坍缩形成的结构（旋转坍缩星）的性质，仅依赖于初始物体的总质量 $m$ 和角动量 $M$，而同它的所有其余特性无关\footnote{为了避免误解，我们提醒读者，这里没有考虑携带净电荷的物体.}.如果物体没有整体旋转（$M=0$），那么坍缩星的外部引力场就是中心对称的史瓦西场\footnote{这一论断受到下述 Israel 定理的强烈支持：在爱因斯坦方程所有无穷远处为伽利略的，且有闭合单叶面 $g_{00}=\mathrm{const},\ t=\mathrm{const}$ 的静态解中，史瓦西解是唯一具有视界（$g_{00}=0$）而在其上没有时空度规奇点的解（证明参见 W. Israel, Phys. Rev. 164, 1776, 1967）.}.

旋转黑洞的引力场由如下轴对称稳态克尔度规给出\footnote{爱因斯坦方程的这个解是 1963 年由克尔以另一形式发现的，后由 R. H. Boyer 和 R. W. Lindquist, 1967 化为 (104.2). 文献中对度规 (104.2) 没有物理概念合格的说明性解析推导，甚至直接核对爱因斯坦方程的这个解也涉及繁杂的计算. 克尔度规是旋转坍缩星唯一的场这个论断，得到一个类似 Israel 定理的定理支持（参见 B. Carter, Phys. Rev. Lett. 26, 331, 1971）.}：
\begin{equation}
\begin{aligned}
  \mathrm{d}s^{2}={}&\left(1-\frac{r_{\mathrm{g}}r}{\rho^{2}}\right)
  \mathrm{d}t^{2}
  -\frac{\rho^{2}}{\varDelta}\,\mathrm{d}r^{2}
  -\rho^{2}\mathrm{d}\theta^{2}\\
  &-\left(r^{2}+a^{2}
  +\frac{r_{\mathrm{g}}ra^{2}}{\rho^{2}}\sin^{2}\theta\right)
  \sin^{2}\theta\,\mathrm{d}\varphi^{2}
  +\frac{2r_{\mathrm{g}}ra}{\rho^{2}}\sin^{2}\theta\,\mathrm{d}\varphi\,\mathrm{d}t,
\end{aligned}
\end{equation}
这里我们已引入符号
\begin{equation}
  \varDelta=r^{2}-r_{\mathrm{g}}r+a^{2},
  \qquad
  \rho^{2}=r^{2}+a^{2}\cos^{2}\theta,
\end{equation}
式中的 $r_{\mathrm{g}}$ 仍然是 $2km/c^{2}$.这个度规依赖于两个恒定参量 $m$ 和 $a$，其意义从度规在远距离 $r$ 处的极限形式看是明显的.精确到 $\sim1/r$ 级的项，我们有：
\begin{equation*}
  g_{00}\approx1-\frac{r_{\mathrm{g}}}{r},\qquad
  g_{03}\approx\frac{r_{\mathrm{g}}a}{r}\sin^{2}\theta;
\end{equation*}

第一式与 (100.18) 比较、第二式与 (104.1) 比较表明，$m$ 是物体的质量，而参量 $a$ 与角动量 $M$ 的关系为
\begin{equation}
  M=ma
\end{equation}
（在通常单位中 $M=mac$）.对于 $a=0$，克尔度规化为标准形式的史瓦西度规 (100.14)\footnote{准确到 $a$ 的一级项，当 $a\ll1$ 时度规 (104.2) 与史瓦西度规只差 $(2r_{\mathrm{g}}a/r)\sin^{2}\theta\,\mathrm{d}\varphi\,\mathrm{d}t$ 这一项，同我们前面关于弱偏离球对称情形的论断一致.}.我们也注意到，形式 (104.2) 明显表现出时间反演对称性：这个变换（$t\to-t$）也改变了转动方向，即角动量的符号（$a\to-a$），所以 $\mathrm{d}s^{2}$ 保持不变.

度规张量 (104.2) 的行列式是
\begin{equation}
  -g=\rho^{4}\sin^{2}\theta.
\end{equation}

我们也给出逆变分量 $g^{ik}$，办法是将它们引入下列四维梯度算符平方的表达式：
\begin{equation}
\begin{aligned}
  g^{ik}\frac{\partial}{\partial x^{i}}
  \frac{\partial}{\partial x^{k}}
  ={}&\frac{1}{\varDelta}\left(r^{2}+a^{2}
  +\frac{r_{\mathrm{g}}ra^{2}}{\rho^{2}}\sin^{2}\theta\right)
  \left(\frac{\partial}{\partial t}\right)^{2}
  -\frac{\varDelta}{\rho^{2}}\left(\frac{\partial}{\partial r}\right)^{2}\\
  &-\frac{1}{\rho^{2}}\left(\frac{\partial}{\partial\theta}\right)^{2}
  -\frac{1}{\varDelta\sin^{2}\theta}
  \left(1-\frac{r_{\mathrm{g}}r}{\rho^{2}}\right)
  \left(\frac{\partial}{\partial\varphi}\right)^{2}
  +\frac{2r_{\mathrm{g}}ra}{\rho^{2}\varDelta}
  \frac{\partial}{\partial\varphi}\frac{\partial}{\partial t}.
\end{aligned}
\end{equation}

当 $m=0$ 时，不存在引力质量，(104.2) 应当变为伽利略度规.实际上，表达式
\begin{equation}
  \mathrm{d}s^{2}=\mathrm{d}t^{2}
  -\frac{\rho^{2}}{r^{2}+a^{2}}\,\mathrm{d}r^{2}
  -\rho^{2}\mathrm{d}\theta^{2}
  -(r^{2}+a^{2})\sin^{2}\theta\,\mathrm{d}\varphi^{2}
\end{equation}

就是用空间扁椭球坐标写出的伽利略度规
\begin{equation*}
  \mathrm{d}s^{2}=\mathrm{d}t^{2}-\mathrm{d}x^{2}-\mathrm{d}y^{2}-\mathrm{d}z^{2},
\end{equation*}

空间扁椭球坐标到笛卡儿坐标的变换由下列公式实现：
\begin{equation*}
\begin{aligned}
  &x=\sqrt{r^{2}+a^{2}}\,\sin\theta\cos\varphi,\\
  &y=\sqrt{r^{2}+a^{2}}\,\sin\theta\sin\varphi,\\
  &z=r\cos\theta;
\end{aligned}
\end{equation*}

曲面 $r=\mathrm{const}$ 是旋转扁椭球：
\begin{equation*}
  \frac{x^{2}+y^{2}}{r^{2}+a^{2}}+\frac{z^{2}}{r^{2}}=1.
\end{equation*}

度规 (104.2) 有一假奇点，正如史瓦西度规 (100.14) 在 $r=r_{\mathrm{g}}$ 有假奇点一样.但是，相较于史瓦西情形在同一个面 $r=r_{\mathrm{g}}$ 上既有 $g_{00}$ 为 0 又有 $g_{11}$ 为无穷大，在克尔度规中，这两个面却是分离的.当 $\rho^{2}=rr_{\mathrm{g}}$ 时，等式 $g_{00}=0$ 成立；这个二次方程的两个根的较大者是
\begin{equation}
  r_0=\frac{r_{\mathrm{g}}}{2}
  +\sqrt{\left(\frac{r_{\mathrm{g}}}{2}\right)^{2}
  -a^{2}\cos^{2}\theta}
  \qquad(g_{00}=0).
\end{equation}

当 $\varDelta=0$ 时，$g_{11}$ 为无穷大；这个方程两个根中较大者是
\begin{equation}
  r_{\mathrm{hor}}=\frac{r_{\mathrm{g}}}{2}
  +\sqrt{\left(\frac{r_{\mathrm{g}}}{2}\right)^{2}-a^{2}}
  \qquad(g_{11}=\infty).
\end{equation}

为简便起见，我们把曲面 $r=r_0$ 记为 $S_0$，$r=r_{\mathrm{hor}}$ 记为 $S_{\mathrm{hor}}$，其物理意义下面再解释.曲面 $S_{\mathrm{hor}}$ 和 $S_0$ 都是旋转扁椭球面\footnote{原著将前者误写为球面.——中译注}；$S_{\mathrm{hor}}$ 包含在 $S_0$ 里，两个面在极点（$\theta=0$ 和 $\theta=\pi$）接触.

正如我们从 (104.8) 和 (104.9) 看到的那样，仅当 $a\leqslant r_{\mathrm{g}}/2$ 时，$S_0$ 和 $S_{\mathrm{hor}}$ 才存在.当 $a>r_{\mathrm{g}}/2$ 时，度规 (104.2) 的性质会发生根本改变，开始显示出物理上不允许的性质——破坏了因果性原理\footnote{从闭合类时世界线的出现就能显示出来这种破坏；这些世界线会使先向过去“运动”然后再朝将来“运动”成为可能. 我们立刻指出，如果把克尔度规延伸到 $S_{\mathrm{hor}}$ 内部，即使 $a<r_{\mathrm{g}}/2$，同样的破坏也会出现. 这表明物理上该度规在 $S_{\mathrm{hor}}$ 内部是不可用的（我们以后将回到这一点）. 由于同样的理由，在二次方程 $g_{00}=0$ 和 $1/g_{11}=0$ 两个较小根所定义的曲面里（它们处于 $S_{\mathrm{hor}}$ 内部），是没有物理意义可言的；参见 B. Carter, Phys. Rev. 147, 1559 (1968).}.

$a>r_{\mathrm{g}}/2$ 时克尔度规失去意义意味着，下列值
\begin{equation}
  a_{\mathrm{max}}=\frac{r_{\mathrm{g}}}{2},
  \qquad
  M_{\mathrm{max}}=\frac{mr_{\mathrm{g}}}{2}
\end{equation}
给出了坍缩星角动量可能的上限.必须把它看成可以任意趋近的极限值，而严格的等式 $a=a_{\mathrm{max}}$ 是不可能的.曲面 $S_0$ 和 $S_{\mathrm{hor}}$ 相应的半径极限值是
\begin{equation}
  r_0=\frac{r_{\mathrm{g}}}{2}(1+\sin\theta),
  \qquad
  r_{\mathrm{hor}}=\frac{r_{\mathrm{g}}}{2}.
\end{equation}

我们将证明曲面 $S_{\mathrm{hor}}$ 是事件视界，运动粒子和光只能从一个方向（朝着内部）穿过它.

先从一般的观点证明，运动粒子世界线的单向穿过性对任何类光超曲面（即每一点的法线都是类光矢量的超曲面）都保持.设超曲面是由方程 $f(x^0,x^1,x^2,x^3)=\mathrm{const}$ 定义的.它的法线沿着四维梯度 $n_i=\partial f/\partial x^i$ 的方向，所以，对于类光超曲面我们有 $n_in^i=0$.换言之，这意味着，法线的方向就躺在该曲面里：沿着该超曲面 $\mathrm{d}f=n_i\,\mathrm{d}x^i=0$，这个方程在四维矢量 $\mathrm{d}x^i$ 和 $n^i$ 的方向一致时满足.由这同样的性质 $n_in^i=0$，超曲面在这同一方向的线元是 $\mathrm{d}s=0$.换言之，沿着这个方向，超曲面在给定点同在该点构造的光锥相切.因此，在类光超曲面上每一点构造的光锥（例如，朝未来方向）完全躺在曲面的一边，并沿自己的母线之一与该超曲面（在这些点）相切.但这正好意味着，粒子或光线（指向未来）的世界线只能往一个方向穿过超曲面.

类光超曲面的这个性质通常在物理上是平凡的：单向穿过这些曲面仅仅表示超光速运动的不可能性（这类例子中最简单的是平直时空中的超曲面 $x=t$）.非平凡的新物理情况出现于类光超曲面并不延伸到空间无限远，以至它的截面 $t=\mathrm{const}$ 是闭合的空间曲面之时；这些曲面是事件视界，其意义与史瓦西球在中心对称引力场中的相同.

克尔场中的曲面 $S_{\mathrm{hor}}$ 是什么呢？对于克尔场中形如 $f(r,\theta)=\mathrm{const}$ 的超曲面，条件 $n_in^i=0$ 具有形式
\begin{equation}
  g^{11}\left(\frac{\partial f}{\partial r}\right)^{2}
  +g^{22}\left(\frac{\partial f}{\partial\theta}\right)^{2}
  =\frac{1}{\rho^{2}}\left[\varDelta\left(\frac{\partial f}{\partial r}\right)^{2}
  +\left(\frac{\partial f}{\partial\theta}\right)^{2}\right]=0
\end{equation}
（$g^{ik}$ 来自 (104.6)）.这个方程在 $S_0$ 上不满足，但在 $S_{\mathrm{hor}}$ 上满足（这里 $\partial f/\partial\theta=0$，$\varDelta=0$）.

把克尔度规延拓到视界曲面以内（像在§102和§103对史瓦西度规做过的那样）是没有物理意义的.这样的延拓会只依赖于和 $S_{\mathrm{hor}}$ 外面的场相同的两个参量（$m$ 和 $a$），由此已经显见，它不可能与坍缩星穿过视界后的命运这个物理问题有关.非球形的效应在共动参考系中根本不能衰减掉，相反，必定会随着物体进一步收缩而增加，所以没有理由预期，视界内部的场会仅仅由物体的质量和角动量来决定\footnote{数学上，这种情况反映了，当我们将克尔度规延伸到 $S_{\mathrm{hor}}$ 以内时（上面提到的）因果性原理的破坏.}.

我们来考虑曲面 $S_0$ 和它与视界之间的空间（克尔场的这个区域称为能层）的性质.

能层的基本性质是，其中没有粒子能够相对于遥远观察者的参考系保持静止：当 $r,\theta,\varphi=\mathrm{const}$ 时，我们有 $\mathrm{d}s^{2}<0$，即间隔是类空的，而粒子的世界线本该类时；变量 $t$ 失去了它的时间特性.因此刚性参考系不能从无穷远延伸入能层，在这个意义上，曲面 $S_0$ 可以称为静界.

粒子在能层内的运动特征与史瓦西场视界之内的情形有本质的不同.在后一种情形下，粒子相对于外部参考系也不能处于静止，不能有 $r=\mathrm{const}$，所有粒子都必须朝中心径向运动.在克尔场的能层中，$\varphi=\mathrm{const}$ 对粒子是不可能的（粒子一定要绕场的对称轴转动），而 $r=\mathrm{const}$ 对于一个粒子是可能的.此外，粒子（和光线）的运动既可以增加也可以减小 $r$，并能从能层出射到外部空间.与最后一点相应，粒子可以从外部区域达到能层：这样的粒子（或光线）到达曲面 $S_0$ 的时间，用遥远观察者的钟 $t$ 测量，对除极点（$S_0$ 和 $S_{\mathrm{hor}}$ 的接触点）之外的所有 $S_0$ 都是有限的.到达这些极点的时间，正如同到达所有 $S_{\mathrm{hor}}$ 点的时间一样，当然还是无限的\footnote{在粒子的能量和角动量取特殊值，使径向速度在 $S_{0}$ 的特殊点为零的特殊情形下，到达这些特殊点的时间也可以是无限的.}.

因为粒子在能层中的转动不可避免，在这个区域写出度规的自然形式是：
\begin{equation}
  \mathrm{d}s^{2}
  =\left(g_{00}-\frac{g_{03}^{2}}{g_{33}}\right)\mathrm{d}t^{2}
  +g_{11}\,\mathrm{d}r^{2}+g_{22}\,\mathrm{d}\theta^{2}
  +g_{33}\left(\mathrm{d}\varphi
  +\frac{g_{03}}{g_{33}}\mathrm{d}t\right)^{2}.
\end{equation}

$\mathrm{d}t^{2}$ 的系数
\begin{equation*}
  g_{00}-\frac{g_{03}^{2}}{g_{33}}
  =\frac{\varDelta}{r^{2}+a^{2}
  +r_{\mathrm{g}}ra^{2}\sin^{2}\theta/\rho^{2}}
\end{equation*}
在 $S_{\mathrm{hor}}$ 外面处处为正（且在 $S_0$ 上不为零）；当 $r=\mathrm{const}$，$\theta=\mathrm{const}$，$\mathrm{d}\varphi=-(g_{03}/g_{33})\,\mathrm{d}t$ 时，间隔 $\mathrm{d}s$ 是类时的.相对于外面的观察者，
\begin{equation}
  -\frac{g_{03}}{g_{33}}
  =\frac{r_{\mathrm{g}}ar}
  {\rho^{2}(r^{2}+a^{2})+r_{\mathrm{g}}ra^{2}\sin^{2}\theta}
\end{equation}
这个量起着广义“能层转动角速度”的作用（其转动方向与中心物体转动的方向一致）\footnote{值得注意的是，对于沿能层边界运动的粒子，固有时间隔并不随 $g_{00}$ 变为零. 在这个意义上，$S_{0}$ 不是“无限红移”面；运动的源（在那里是不能静止的）送出的光信号频率，在遥远观察者看来并不变为零. 我们记得，在中心对称场中，史瓦西球上既不可能有静止的源，也不可能有运动的源（因为类光曲面不能包含类时世界线）. 在那里产生“无限红移”是因为，随着 $r\to r_{\mathrm{g}}$，固有时间隔 $\mathrm{d}\tau=\sqrt{g_{00}}\,\mathrm{d}t$（对于给定的 $\mathrm{d}t$），由相对于参考系静止的钟测量，会变为零.}.

粒子的能量，定义为作用量 $S$ 对于粒子（沿轨道同步的）固有时 $\tau$ 的导数 $-\partial S/\partial\tau$，它恒为正值（参见§88）.但是，正如§88中说明过的那样，粒子在与时间 $t$ 无关的场中运动期间，定义为 $-\partial S/\partial t$ 的能量 $\mathcal{E}_0$ 是守恒的；这个量与四维动量 $p_0=mu_0=mg_{0i}\,\mathrm{d}x^i/\mathrm{d}s$ 的协变分量一致（$m$ 是粒子的质量）.变量 $t$（按遥远观察者的钟测量的时间）在能层内不再具有时间特性，这个事实产生了如下特殊的情况：在这个区间内 $g_{00}<0$，因此
\begin{equation*}
  \mathcal{E}_0=m(g_{00}u^{0}+g_{03}u^{3})
  =m\left(g_{00}\frac{\mathrm{d}t}{\mathrm{d}s}
  +g_{03}\frac{\mathrm{d}\varphi}{\mathrm{d}s}\right)
\end{equation*}
这个量能够为负值.因为在外部空间（那里 $t$ 为时间）中，能量 $\mathcal{E}_0$ 不能为负，所以 $\mathcal{E}_0<0$ 的粒子不能从外部落入能层.出现这种粒子的可能来源，是进入能层的物体分裂为（例如说）两部分，一部分被俘获进入“负能”轨道，这部分不再能从能层出射，最后被俘获进视界之内.第二部分可以返回外部空间；因为 $\mathcal{E}_0$ 是守恒的加性量，这部分的能量大于初始物体的能量，所以我们得到了从转动黑洞抽取出来的能量 (R. Penrose, 1969).

最后，我们注意到，虽然 $S_0$ 面对于时空度规并不奇异，但纯空间度规（在参考系 (104.2) 中）在那里却有奇点.在 $S_0$ 以外（那里变量 $t$ 具有时间特性），空间度规张量根据 (84.7) 计算，而空间距离元具有形式
\begin{equation}
  \mathrm{d}l^{2}=\frac{\rho^{2}}{\varDelta}\,\mathrm{d}r^{2}
  +\rho^{2}\mathrm{d}\theta^{2}
  +\frac{\varDelta\sin^{2}\theta}{1-rr_{\mathrm{g}}/\rho^{2}}\,\mathrm{d}\varphi^{2}.
\end{equation}

在 $S_0$ 附近，纬线长度（$\theta=\mathrm{const}$，$r=\mathrm{const}$）按规律 $2\pi a\sin^{2}\theta/\sqrt{g_{00}}$ 趋于无穷大.当时钟沿此闭合曲线同步时，其读数之差（参见 (88.5)）在这里也趋于无穷大.

\begin{xiti}

\noindent{\heiti 习题1}\quad 对于在克尔场中运动的粒子，在哈密顿-雅可比方程中进行变量分离 (B. Carter, 1968).

解：在哈密顿-雅可比方程
\begin{equation*}
  g^{ik}\frac{\partial S}{\partial x^{i}}
  \frac{\partial S}{\partial x^{k}}-m^{2}=0
\end{equation*}
中（$m$ 是粒子的质量，不要与中心物体的质量混淆），度规 $g^{ik}$ 来自 (104.6)，时间 $t$ 和角 $\varphi$ 是循环变量；因此它们以 $-\mathcal{E}_0t+L\varphi$ 的形式进入作用量 $S$，这里 $\mathcal{E}_0$ 是守恒的能量，$L$ 表示角动量沿场对称轴的分量.显然，变量 $r$ 和 $\theta$ 也可以分离.把 $S$ 写为形式
\begin{equation}
  S=-\mathcal{E}_0t+L\varphi+S_r(r)+S_\theta(\theta),
  \tag{1}
\end{equation}
我们把哈密顿-雅可比方程化为两个常微分方程（参见本教程第一卷§48）：
\begin{equation*}
  \left(\frac{\mathrm{d}S_\theta}{\mathrm{d}\theta}\right)^{2}
  +\left(a\mathcal{E}_0\sin\theta-\frac{L}{\sin\theta}\right)^{2}
  +a^{2}m^{2}\cos^{2}\theta=K,
\end{equation*}

\begin{equation}
  \left(\frac{\mathrm{d}S_r}{\mathrm{d}r}\right)^{2}
  -\frac{1}{\varDelta}\left[(r^{2}+a^{2})\mathcal{E}_0-aL\right]^{2}
  +m^{2}r^{2}=-K,
  \tag{2}
\end{equation}
式中 $K$（分离参数）是一个新的任意常数.于是函数 $S_\theta$ 和 $S_r$ 可以通过简单求积分决定.

粒子的四维动量是
\begin{equation*}
  p^{i}=m\frac{\mathrm{d}x^{i}}{\mathrm{d}s}
  =g^{ik}p_k=-g^{ik}\frac{\partial S}{\partial x^{k}}.
\end{equation*}

用 (1) 和 (2) 计算这个方程的右边，我们得到下列方程：
\begin{equation}
  m\frac{\mathrm{d}t}{\mathrm{d}s}
  =-\frac{r_{\mathrm{g}}ra}{\rho^{2}\varDelta}L
  +\frac{\mathcal{E}_0}{\varDelta}
  \left(r^{2}+a^{2}
  +\frac{r_{\mathrm{g}}ra^{2}}{\rho^{2}}\sin^{2}\theta\right),
  \tag{3}
\end{equation}

\begin{equation}
  m\frac{\mathrm{d}\varphi}{\mathrm{d}s}
  =\frac{L}{\varDelta\sin^{2}\theta}
  \left(1-\frac{r_{\mathrm{g}}r}{\rho^{2}}\right)
  +\frac{r_{\mathrm{g}}ra}{\rho^{2}\varDelta}\mathcal{E}_0,
  \tag{4}
\end{equation}

\begin{equation}
  m^{2}\left(\frac{\mathrm{d}r}{\mathrm{d}s}\right)^{2}
  =\frac{1}{\rho^{4}}
  \left[(r^{2}+a^{2})\mathcal{E}_0-aL\right]^{2}
  -\frac{\varDelta}{\rho^{4}}(K+m^{2}r^{2}),
  \tag{5}
\end{equation}

\begin{equation}
  m^{2}\left(\frac{\mathrm{d}\theta}{\mathrm{d}s}\right)^{2}
  =\frac{1}{\rho^{4}}(K-a^{2}m^{2}\cos^{2}\theta)
  -\frac{1}{\rho^{4}}
  \left(a\mathcal{E}_0\sin\theta-\frac{L}{\sin\theta}\right)^{2}.
  \tag{6}
\end{equation}

这些积分是运动方程（测地方程）的初积分.轨道方程和沿轨道坐标对时间的依赖关系可以从 (3)—(6) 或者直接从方程
\begin{equation*}
  \frac{\partial S}{\partial\mathcal{E}_0}=\mathrm{const},\qquad
  \frac{\partial S}{\partial L}=\mathrm{const},\qquad
  \frac{\partial S}{\partial K}=\mathrm{const}
\end{equation*}
求得.

对于光线的情况，我们必须在方程 (3)—(6) 右边令 $m=0$，并用 $\omega_0$ 代替 $\mathcal{E}_0$（参见§101），而左边的导数 $\mathrm{d}m/\mathrm{d}s$ 必须代之以对（沿光线变化的）参数 $\lambda$ 的导数 $\mathrm{d}/\mathrm{d}\lambda$（参见§87末）.

如从对称性论据已经清楚的那样，方程 (4)—(6) 只容许沿物体转动轴的纯径向运动.由与此相同的考虑显见，运动可能在一个“平面”上进行的必要条件是该平面是赤道面.在该情形下，令 $\theta=\pi/2$，从条件 $\mathrm{d}\theta/\mathrm{d}s=0$ 用 $\mathcal{E}_0$ 和 $L$ 表示 $K$，我们得到如下形式的运动方程：
\begin{equation}
  m\frac{\mathrm{d}t}{\mathrm{d}s}
  =-\frac{r_{\mathrm{g}}a}{r\varDelta}L
  +\frac{\mathcal{E}_0}{\varDelta}
  \left(r^{2}+a^{2}+\frac{r_{\mathrm{g}}a^{2}}{r}\right)
  \tag{7}
\end{equation}

\begin{equation}
  m\frac{\mathrm{d}\varphi}{\mathrm{d}s}
  =\frac{M}{\varDelta}\left(1-\frac{r_{\mathrm{g}}}{r}\right)
  +\frac{r_{\mathrm{g}}a}{r\varDelta}\mathcal{E}_0,
  \tag{8}
\end{equation}

\begin{equation}
  m^{2}\left(\frac{\mathrm{d}r}{\mathrm{d}s}\right)^{2}
  =\frac{1}{r^{4}}
  \left[(r^{2}+a^{2})\mathcal{E}_0-aL\right]^{2}
  -\frac{\varDelta}{r^{4}}
  \left[(a\mathcal{E}_0-L)^{2}+m^{2}r^{2}\right]
  \tag{9}
\end{equation}

\noindent{\heiti 习题2}\quad 对于在极限克尔场（$a\to r_{\mathrm{g}}/2$）赤道面内运动的粒子，求最接近中心的稳定圆轨道的半径 (R. Ruffini, J. A. Wheeler, 1969).

解：同§102习题1一样处理，引入由
\begin{equation*}
  \left[(r^{2}+a^{2})U(r)-aL\right]^{2}
  -\varDelta\left[(aU(r)-L)^{2}+r^{2}m^{2}\right]=0
\end{equation*}
定义的有效势能 $U(r)$（对于 $\mathcal{E}_0=U$，方程 (9) 的右边变为零）.稳定轨道的半径由函数 $U(r)$ 的极小值，即方程 $U(r)=\mathcal{E}_0,\ U'(r)=0$ 在 $U''(r)>0$ 条件下的联立解确定.最接近中心的轨道相应于 $U''(r_{\mathrm{min}})=0$；对于 $r<r_{\mathrm{min}}$，函数 $U(r)$ 没有极小值.结果我们得到如下的运动参数值：

(a) 对 $L<0$（运动与坍缩星转动方向相反）
\begin{equation*}
  \frac{r_{\mathrm{min}}}{r_{\mathrm{g}}}=\frac{9}{2},
  \qquad
  \frac{\mathcal{E}_0}{m}=\frac{5}{3\sqrt{3}},
  \qquad
  \frac{L}{mr_{\mathrm{g}}}=\frac{11}{3\sqrt{3}}.
\end{equation*}

(b) 对 $L>0$（运动沿着坍缩星转动方向），当 $a\to\dfrac{r_{\mathrm{g}}}{2}$ 时半径 $r_{\mathrm{min}}$ 趋向视界半径.令 $a=\dfrac{r_{\mathrm{g}}}{2}(1+\delta)$，当 $\delta\to0$ 时我们得到：
\begin{equation*}
  \frac{r_{\mathrm{hor}}}{r_{\mathrm{g}}}
  =\frac{1}{2}(1+\sqrt{2\delta}),\qquad
  \frac{r_{\mathrm{min}}}{r_{\mathrm{g}}}
  =\frac{1}{2}\left[1+(4\delta)^{1/3}\right].
\end{equation*}
于是
\begin{equation*}
  \frac{\mathcal{E}_0}{m}
  =\frac{L}{mr_{\mathrm{g}}}
  =\frac{1}{\sqrt{3}}\left[1+(4\delta)^{1/3}\right].
\end{equation*}

我们注意到，$r_{\mathrm{min}}/r_{\mathrm{hor}}$ 总是大于 1，即轨道不会处于视界以内.这本该如此：视界是类光超曲面，运动粒子的类时世界线是不可能置身其上的.

\end{xiti}

\section{物体远距离处的引力场}

我们来考虑远离场源物体处的稳态引力场，并决定其按 $1/r$ 幂展开式的前几项.

离物体远处的场很弱.这意味着那里的时空度规几乎是伽利略度规，即我们可以选择一个参考系，其中度规张量的分量几乎等于它们的伽利略值：
\begin{equation}
  g_{00}^{(0)}=1,\qquad g_{0\alpha}^{(0)}=0,\qquad
  g_{\alpha\beta}^{(0)}=-\delta_{\alpha\beta}.
\end{equation}

因而我们可以把 $g_{ik}$ 写成如下形式：
\begin{equation}
  g_{ik}=g_{ik}^{(0)}+h_{ik},
\end{equation}
式中 $h_{ik}$ 是由引力场确定的小修正.

在对张量 $h_{ik}$ 的运算中，我们约定用“未扰动的”度规来升和降它们的指标：$h_i^k=g^{(0)kl}h_{il}$，等等.这里我们必须把 $h^{ik}$ 同度规张量 $g^{ik}$ 的逆变分量的修正区分开.后者由解如下方程确定：
\begin{equation*}
  g_{il}g^{lk}=(g_{il}^{(0)}+h_{il})g^{lk}=\delta_i^k;
\end{equation*}
所以，精确到二阶项，我们得到：
\begin{equation}
  g^{ik}=g^{ik(0)}-h^{ik}+h_l^ih^{lk}.
\end{equation}

精确到同样精度，度规张量的行列式是
\begin{equation}
  g=g^{(0)}\left(1+h+\frac{1}{2}h^{2}
  -\frac{1}{2}h_k^ih_i^k\right),
\end{equation}
式中 $h\equiv h_i^i$.

我们现在就要强调，$h_{ik}$ 很小这个条件绝对没有对参考系作唯一确定的选择.如果这个条件对任一参考系满足，在作任意变换 $x'^i=x^i+\xi^i$ 后它也将满足，这里 $\xi^i$ 是些小量.按照 (94.3)，张量 $h_{ik}$ 就将变为
\begin{equation}
  h'_{ik}=h_{ik}
  -\frac{\partial\xi_i}{\partial x^{k}}
  -\frac{\partial\xi_k}{\partial x^{i}},
\end{equation}
这里 $\xi_i=g_{ik}^{(0)}\xi^k$（因为 $g_{ik}^{(0)}$ 是常数，(94.3) 中的协变导数在目前情形中化为普通导数）\footnote{对于稳态场，自然只容许那些不破坏 $g_{ik}$ 的时间无关性的变换，即 $\xi^{i}$ 必须只是空间坐标的函数.}.

在一阶近似下，精确到 $1/r$ 的项，对伽利略值的小修正由中心对称史瓦西度规展开式中相应的项给出.由于上面提到的参考系（无穷远处为伽利略参考系）选择的不定性，$h_{ik}$ 的具体形式依赖于如何定义径向坐标 $r$.因此，如果史瓦西度规写成 (100.14) 的形式，其大 $r$ 处展开式中的前几项就由表达式 (100.18) 给出.从空间球坐标变换到笛卡儿坐标（为此我们必须写出 $\mathrm{d}r=n_{\alpha}\,\mathrm{d}x^{\alpha}$，这里 $n$ 是 $r$ 方向的单位矢量），我们得到下列值：
\begin{equation}
  h_{00}^{(1)}=-\frac{r_{\mathrm{g}}}{r},\qquad
  h_{\alpha\beta}^{(1)}=-\frac{r_{\mathrm{g}}}{r}n_{\alpha}n_{\beta},
  \qquad
  h_{0\alpha}^{(1)}=0,
\end{equation}
式中 $r_{\mathrm{g}}=2km/c^{2}$.\footnote{然而如果我们从有各向同性空间坐标的史瓦西度规出发（参见§100习题4），我们会得到：$h_{00}^{(1)}=-r_{\mathrm{g}}/r,\ h_{\alpha\beta}^{(1)}=-(r_{\mathrm{g}}/r)\delta_{\alpha\beta},\ h_{0\alpha}^{(1)}=0$ (105.6a). 从 (105.6) 移到 (105.6a) 是通过变换 (105.5) 实现的，其中令 $\xi^{0}=0,\ \xi^{\alpha}=-r_{\mathrm{g}}x^{\alpha}/2r$.}

在正比于 $1/r^{2}$ 的二阶项中，有两种不同起源的项.有些项来自爱因斯坦方程对一阶项的非线性效应.因为后者只依赖于质量（而不依赖于物体的其他特性），故这种二阶项也只依赖于质量.因而很清楚，这些项可以通过展开史瓦西度规得到.在这些坐标中，我们有
\begin{equation}
  h_{00}^{(2)}=0,\qquad
  h_{\alpha\beta}^{(2)}
  =-\left(\frac{r_{\mathrm{g}}}{r}\right)^{2}n_{\alpha}n_{\beta}.
\end{equation}

其余的二阶项来自线性化场方程的相应解.考虑到后面的应用，我们将超出这里稳定场的需要，用形式上更一般的公式来进行方程的线性化；开始时我们不用场的稳定性.

对于小的 $h_{ik}$，用其导数表示的量 $\Gamma_{kl}^i$ 也很小.忽略高于一次的幂，在曲率张量 (92.1) 中我们可以只保留第一个括号中的项：
\begin{equation}
  R_{iklm}=\frac{1}{2}\left(
  \frac{\partial^{2}h_{im}}{\partial x^{k}\partial x^{l}}
  +\frac{\partial^{2}h_{kl}}{\partial x^{i}\partial x^{m}}
  -\frac{\partial^{2}h_{km}}{\partial x^{i}\partial x^{l}}
  -\frac{\partial^{2}h_{il}}{\partial x^{k}\partial x^{m}}\right).
\end{equation}

对于里奇张量，准确到同样精度，我们有：
\begin{equation*}
  R_{ik}=g^{lm}R_{limk}\approx g^{lm(0)}R_{limk},
\end{equation*}
或者
\begin{equation}
  R_{ik}=\frac{1}{2}\left(
  -g^{lm(0)}\frac{\partial^{2}h_{ik}}{\partial x^{l}\partial x^{m}}
  +\frac{\partial^{2}h_i^l}{\partial x^{k}\partial x^{l}}
  +\frac{\partial^{2}h_k^l}{\partial x^{i}\partial x^{l}}
  -\frac{\partial^{2}h}{\partial x^{i}\partial x^{k}}\right).
\end{equation}

表达式 (105.9) 可以通过使用参考系选择中留下的任意性得到简化.我们可以施予 $h_{ik}$ 四个（任意函数 $\xi^i$ 的数目）附加条件
\begin{equation}
  \frac{\partial\psi_i^k}{\partial x^{k}}=0,
  \qquad
  \psi_i^k=h_i^k-\frac{1}{2}\delta_i^kh.
\end{equation}

于是 (105.9) 中的最后三项彼此相消，留下
\begin{equation}
  R_{ik}=-\frac{1}{2}g^{lm(0)}
  \frac{\partial^{2}h_{ik}}{\partial x^{l}\partial x^{m}}.
\end{equation}

在我们这里考虑的稳定情形下，$h_{ik}$ 不依赖于时间，表达式 (105.11) 化为 $R_{ik}=\Delta h_{ik}/2$，式中 $\Delta$ 是三维空间坐标中的拉普拉斯算符.因此真空中的爱因斯坦场方程化为拉普拉斯方程
\begin{equation}
  \Delta h_{ik}=0,
\end{equation}
连同附加条件 (105.10)，后者取形式
\begin{equation}
  \frac{\partial}{\partial x^{\beta}}
  \left(h_{\alpha}^{\beta}-\frac{1}{2}h\delta_{\alpha}^{\beta}\right)=0,
\end{equation}

\begin{equation}
  \frac{\partial}{\partial x^{\beta}}h_0^{\beta}=0.
\end{equation}

我们注意到，这些条件仍然没有完全确定参考系的唯一选择.显而易见，如果 $h_{ik}$ 满足方程 (105.13)—(105.14)，则同样的条件也将被 (105.5) 的 $h'_{ik}$ 满足，只要 $\xi^i$ 满足方程
\begin{equation}
  \Delta\xi^i=0.
\end{equation}

分量 $h_{00}$ 必定由三维拉普拉斯方程的标量解给出.我们知道，正比于 $1/r^{2}$ 的这样一个解具有形式 $\boldsymbol{a}\cdot\nabla(1/r)$，这里 $\boldsymbol{a}$ 是一个常矢量.但 $h_{00}$ 中这种类型的项总能够通过简单地在 $1/r$ 的一阶项中移动坐标原点消去.因此，这样一项的存在只是表明坐标原点选择得不好，并没有什么意义.

分量 $h_{0\alpha}$ 由拉普拉斯方程的矢量解给出，即它们必须具有形式
\begin{equation*}
  h_{0\alpha}=\lambda_{\alpha\beta}
  \frac{\partial}{\partial x^{\beta}}\frac{1}{r},
\end{equation*}
式中 $\lambda_{\alpha\beta}$ 是一个常张量.条件 (105.14) 给出
\begin{equation*}
  \lambda_{\alpha\beta}
  \frac{\partial^{2}}{\partial x^{\alpha}\partial x^{\beta}}
  \frac{1}{r}=0,
\end{equation*}
由此可知，$\lambda_{\alpha\beta}$ 必须具有形式 $a_{\alpha\beta}+\lambda\delta_{\alpha\beta}$，这里 $a_{\alpha\beta}$ 是一个反对称张量.但是形如 $\lambda\dfrac{\partial}{\partial x^{\alpha}}\dfrac{1}{r}$ 的解可以通过变换 (105.5) 并令 $\xi^0=\lambda/r,\ \xi^{\alpha}=0$（满足条件 (105.15)）消去.因此，唯一具有真实意义的解是
\begin{equation*}
  h_{0\alpha}=a_{\alpha\beta}
  \frac{\partial}{\partial x^{\beta}}\frac{1}{r}.
\end{equation*}

最后，通过类似但更为繁杂的论证可以证明，通过适当的空间坐标变换，总能消去拉普拉斯方程张量解（对 $\alpha$ 和 $\beta$ 对称）给出的量 $h_{\alpha\beta}$.

至于张量 $a_{\alpha\beta}$，它同总角动量张量 $M_{\alpha\beta}$ 有关，$h_{0\alpha}$ 的最终表达式具有形式
\begin{equation}
  h_{0\alpha}^{(2)}
  =\frac{2k}{c^{3}}M_{\alpha\beta}
  \frac{\partial}{\partial x^{\beta}}\frac{1}{r}
  =-\frac{2k}{c^{3}}M_{\alpha\beta}\frac{n_{\beta}}{r^{2}}.
\end{equation}

我们通过计算积分 (96.17) 来证明这一点.

角动量 $M_{\alpha\beta}$ 只同 $h_{0\alpha}$ 有关，所以在计算它的时候，我们可以假设所有其他的分量 $h_{ik}$ 都不存在.准确到 $h_{0\alpha}$ 中的二阶项，从 (96.2)—(96.3) 我们有（注意 $g^{\alpha0}=-h^{\alpha0}=h_{\alpha0}$，而 $-g$ 与 1 只差二阶项）：
\begin{equation*}
  h^{\alpha0\beta}
  =\frac{c^{4}}{16\pi k}\frac{\partial}{\partial x^{\gamma}}
  (g^{\alpha0}g^{\beta\gamma}-g^{\gamma0}g^{\alpha\beta})
  =-\frac{c^{4}}{16\pi k}\frac{\partial}{\partial x^{\gamma}}
  (h_{\alpha0}\delta_{\beta\gamma}-h_{\gamma0}\delta_{\alpha\beta}).
\end{equation*}

将 (105.16) 代入这里，导数符号下的第二项变为零，而第一项给出
\begin{equation*}
  h^{\alpha0\beta}
  =-\frac{c}{8\pi}M_{\alpha\gamma}
  \frac{\partial^{2}}{\partial x^{\beta}\partial x^{\gamma}}\frac{1}{r}
  =-\frac{c}{8\pi}M_{\alpha\gamma}
  \frac{3n_{\beta}n_{\gamma}-\delta_{\beta\gamma}}{r^{3}}.
\end{equation*}

用这个表达式在半径 $r$ 的球面上进行 (96.16) 中的积分（$\mathrm{d}f_\gamma=n_\gamma r^2\mathrm{d}o$），我们得到：
\begin{equation*}
\begin{aligned}
  \frac{1}{c}\oint(x^{\alpha}h^{\beta0\gamma}
  -x^{\beta}h^{\alpha0\gamma})\,\mathrm{d}f_{\gamma}
  &=-\frac{1}{4\pi}\int(n_{\alpha}n_{\gamma}M_{\beta\gamma}
  -n_{\beta}n_{\gamma}M_{\alpha\gamma})\,\mathrm{d}o\\
  &=-\frac{1}{3}(\delta_{\alpha\gamma}M_{\beta\gamma}
  -\delta_{\beta\gamma}M_{\alpha\gamma})
  =\frac{2}{3}M_{\alpha\beta}.
\end{aligned}
\end{equation*}

类似的计算给出：
\begin{equation*}
  \frac{1}{c}\oint\lambda^{\alpha0\gamma\beta}\,\mathrm{d}f_{\gamma}
  =-\frac{c^{3}}{16\pi k}\oint
  (h_{\alpha0}\,\mathrm{d}f_{\beta}-h_{\beta0}\,\mathrm{d}f_{\alpha})
  =\frac{1}{3}M_{\alpha\beta}.
\end{equation*}

把这两个量相加，我们就得到要求的 $M_{\alpha\beta}$ 值.

我们强调指出，在一般情形下，当物体附近的场可能并不很弱时，$M_{\alpha\beta}$ 是物体及其引力场的总角动量.仅当场在所有距离都很弱时，才能忽略它对角动量的贡献\footnote{如果转动物体是球形的，$\boldsymbol{M}$ 的方向是物体外整个空间的场唯一可分辨的方向. 如果场处处（而不只是离物体远处）很弱，公式 (105.16) 在物体外面的整个空间都有效. 在场的中心对称部分不是处处很弱，但球形物体转动得足够慢的情形下，这个公式在整个空间仍然成立（参见习题1）.}.

公式 (105.6)—(105.7) 和 (105.16) 解决了我们的问题，准确到 $1/r^{2}$ 的项\footnote{变换 (105.5) 对于 $\xi^{0}=0,\ \xi^{\alpha}=\xi^{\alpha}(x^{1},x^{2},x^{3})$ 并不改变 $h_{0\alpha}$. 因此，表达式 (105.16) 不依赖于坐标 $r$ 的选择.}.度规张量的协变分量是：
\begin{equation}
  g_{ik}=g_{ik}^{(0)}+h_{ik}^{(1)}+h_{ik}^{(2)}.
\end{equation}

按照 (105.3)，达到与此相同的精度的逆变分量是
\begin{equation}
  g^{ik}=g^{ik(0)}-h^{ik(1)}-h^{ik(2)}
  +h_l^{i(1)}h^{lk(1)}.
\end{equation}

公式 (105.16) 可以用矢量形式重新写为\footnote{准确到假设的精度，矢量 $g_{\alpha}=-g_{0\alpha}/g_{00}\approx-g_{0\alpha}$. 与此同理，在定义矢积和旋度时（参见§88习题1的脚注）我们必须令 $\gamma=1$，所以可以按对笛卡儿矢量常用的意义来理解它们.}：
\begin{equation}
  \boldsymbol{g}=\frac{2k}{c^{3}r^{2}}\boldsymbol{n}\times\boldsymbol{M}.
\end{equation}
式中 $\boldsymbol{M}$ 是物体的总角动量矢量.我们在§88习题1中曾经表明，在稳态引力场中，有一“科里奥利力”作用于物体上，它与物体在以角速度
\begin{equation*}
  \Omega=\frac{c}{2}\sqrt{g_{00}}\operatorname{rot}\boldsymbol{g}
\end{equation*}
旋转的参考系中所受的力相同.因而我们可以说，在旋转物体的场中，作用于远处粒子上的科里奥利力的强度相应于角速度：
\begin{equation}
  \Omega\approx\frac{c}{2}\operatorname{rot}\boldsymbol{g}
  =\frac{k}{c^{2}r^{3}}\left[\boldsymbol{M}
  -3\boldsymbol{n}(\boldsymbol{M}\cdot\boldsymbol{n})\right].
\end{equation}

最后，我们依照积分 (96.16)，用表达式 (105.6) 来计算引力物体的总能量.从公式 (96.2)—(96.3) 计算 $h^{ikl}$ 的必要分量，在需要的精度下（保留 $1/r^{2}$ 阶的项）我们得到：
\begin{equation*}
  h^{00\alpha}
  =\frac{c^{4}}{16\pi k}\frac{\partial}{\partial x^{\beta}}
  (g^{00}g^{\alpha\beta})
  =\frac{mc^{2}}{8\pi}\frac{\partial}{\partial x^{\beta}}
  \left(-\frac{\delta^{\alpha\beta}}{r}
  +\frac{x^{\alpha}x^{\beta}}{r^{3}}\right)
  =\frac{mc^{2}}{4\pi}\frac{n^{\alpha}}{r^{2}}.
\end{equation*}

在半径为 $r$ 的球上作 (96.16) 中的积分，最后得
\begin{equation}
  P^{\alpha}=0,\qquad P^{0}=mc,
\end{equation}

这自然是个预期的结果.它表示了物体的“引力”质量和“惯性”质量相等的事实（“引力”质量是决定物体产生的引力场的质量，这就是出现在引力场的度规张量中，或者特殊情形下，牛顿定律中的质量；“惯性”质量是决定物体的能量与动量之比的质量；特别地，物体的静能等于这个质量乘以 $c^{2}$）.

在恒定引力场情形下，可以推出一个物质加场总能量的简单表达式，它在形式上只对物质所占的空间进行积分.为了做到这一点我们可以，比方说，从如下表达式（它在所有量与 $x^0$ 无关时成立）出发\footnote{从 (92.7) 我们有 $R_{0}^{0}=g^{0i}R_{i0}=g^{0i}\left(\dfrac{\partial\Gamma_{i0}^{l}}{\partial x^{l}}+\Gamma_{i0}^{l}\Gamma_{lm}^{m}-\Gamma_{il}^{m}\Gamma_{0m}^{l}\right)$，用 (86.5) 和 (86.8) 我们发现，这个表达式可以写为 $R_{0}^{0}=\dfrac{1}{\sqrt{-g}}\dfrac{\partial}{\partial x^{i}}\left(\sqrt{-g}\,g^{0i}\Gamma_{i0}^{l}\right)+g^{im}\Gamma_{ml}^{0}\Gamma_{i0}^{l}$；用相同关系 (86.8) 容易证明，右边第二项恒等于 $-\dfrac{1}{2}\Gamma_{lm}^{0}\dfrac{\partial g^{lm}}{\partial x^{0}}$，因所有的量都与 $x^{0}$ 无关，故等于零. 最后将第一项中对 $l$ 的求和换成对 $\alpha$ 求和，我们就得到 (105.22).}：
\begin{equation}
  R_0^0=\frac{1}{\sqrt{-g}}
  \frac{\partial}{\partial x^{\alpha}}
  \left(\sqrt{-g}\,g^{i0}\Gamma_{0i}^{\alpha}\right).
\end{equation}

在（三维）空间中对 $R_0^0\sqrt{-g}$ 作积分，用三维高斯公式得到：
\begin{equation*}
  \int R_0^0\sqrt{-g}\,\mathrm{d}V
  =\oint\sqrt{-g}\,g^{i0}\Gamma_{0i}^{\alpha}\,\mathrm{d}f_{\alpha}.
\end{equation*}

取足够远的表面作这个积分，对 $g_{ik}$ 用表达式 (105.6)，经简单的计算后得到：
\begin{equation*}
  \int R_0^0\sqrt{-g}\,\mathrm{d}V
  =\frac{4\pi k}{c^{2}}m
  =\frac{4\pi k}{c^{3}}P^{0}.
\end{equation*}

再注意到，按照场方程，
\begin{equation*}
  R_0^0=\frac{8\pi k}{c^{4}}\left(T_0^0-\frac{1}{2}T\right)
  =\frac{4\pi k}{c^{4}}(T_0^0-T_1^1-T_2^2-T_3^3),
\end{equation*}
我们得到要求的公式：
\begin{equation}
  P^{0}=mc
  =\frac{1}{c}\int(T_0^0-T_1^1-T_2^2-T_3^3)\sqrt{-g}\,\mathrm{d}V.
\end{equation}

这个公式仅用物质的能量动量张量表达出物质和恒定引力场的总能量（即物体的总质量）(R. Tolman, 1930).我们记得，在中心对称场情形下，这个量还有另一个表达式——公式 (100.23).

\begin{xiti}

\noindent{\heiti 习题1}\quad 证明在慢转动（$M\ll cmr_{\mathrm{g}}$）但不要求场的中心对称部分很小的条件下，公式 (105.16) 对于旋转椭球体外部整个空间的场保持有效 (A. G. Doroshkevich, Ya. B. Zel'dovich, I. D. Novikov, 1965; V. Gurovich, 1965).

解：在空间球坐标中（$x^1=r,\ x^2=\theta,\ x^3=\varphi$），公式 (105.16) 写为：
\begin{equation}
  h_{03}=\frac{2kM}{rc^{2}}\sin^{2}\theta.
  \tag{1}
\end{equation}

把这个量看成是对史瓦西度规 (100.14) 的小修正，我们必须验证按 $h_{03}$ 线性化的方程 $R_{03}=0$ 得以满足（因为在其他场方程中修正项恒为零）.$R_{03}$ 可用§95习题中的公式 (4) 计算，此处线性化意味着，三维张量运算应该用“未扰动”的度规 (100.15) 进行.结果我们得到如下方程
\begin{equation*}
  \left(1-\frac{r_{\mathrm{g}}}{r}\right)
  \frac{\partial^{2}h_{03}}{\partial r^{2}}
  +\frac{2r_{\mathrm{g}}}{r^{3}}h_{03}
  +\frac{\sin\theta}{r^{2}}\frac{\partial}{\partial\theta}
  \left(\frac{1}{\sin\theta}\frac{\partial h_{03}}{\partial\theta}\right)=0,
\end{equation*}
表达式 (1) 确实满足它.

\noindent{\heiti 习题2}\quad 求在转动的中心物体的场中运动粒子轨道的系统（“长期”）移动 (J. Lense, H. Thirring, 1918).

解：由于所有的相对论效应都很小，它们彼此线性地叠加，所以在计算由中心物体转动所产生的效应时，我们可以忽略§101中考虑过的非牛顿的中心对称力场；换言之，在进行计算时我们可以假设，所有的 $h_{ik}$ 中只有 $h_{0\alpha}$ 不等于零.

粒子经典轨道的取向决定于两个守恒量：粒子的轨道角动量 $\boldsymbol{M}=\boldsymbol{r}\times\boldsymbol{p}$ 和矢量
\begin{equation*}
  \boldsymbol{A}
  =\frac{\boldsymbol{p}}{m}\times\boldsymbol{M}
  -\frac{kmm'\boldsymbol{r}}{r},
\end{equation*}
后者的守恒是专对牛顿场 $\varphi=-km'/r$ 而言的（式中 $m'$ 是中心物体的质量）.参见本教程第一卷§15.矢量 $\boldsymbol{M}$ 垂直于轨道平面，而矢量 $\boldsymbol{A}$ 沿椭圆的长轴指向近日点（且其大小等于 $kmm'e$，这里 $e$ 是轨道的偏心率）.要求的轨道的长期移动可以用这些矢量方向的改变来描述.

在场 (105.19) 中运动粒子的拉格朗日函数是
\begin{equation}
  L=-mc\frac{\mathrm{d}s}{\mathrm{d}t}=L_0+\delta L,
  \qquad
  \delta L=mc\,\boldsymbol{g}\cdot\boldsymbol{v}
  =\frac{2km}{c^{2}r^{3}}\boldsymbol{M}'\cdot\boldsymbol{v}\times\boldsymbol{r}
  \tag{1}
\end{equation}
（这里我们将中心物体的角动量记为 $\boldsymbol{M}'$ 以区别于粒子的角动量 $\boldsymbol{M}$）.于是，哈密顿函数是（参见本教程第一卷 (40.7)）：
\begin{equation*}
  \mathcal{H}=\mathcal{H}_0+\delta\mathcal{H},
  \qquad
  \delta\mathcal{H}
  =\frac{2k}{c^{2}r^{3}}\boldsymbol{M}'\cdot\boldsymbol{r}\times\boldsymbol{p}.
\end{equation*}

用哈密顿方程 $\dot{\boldsymbol{r}}=\partial\mathcal{H}/\partial\boldsymbol{p}$，$\dot{\boldsymbol{p}}=-\partial\mathcal{H}/\partial\boldsymbol{r}$ 来计算导数 $\dot{\boldsymbol{M}}=\dot{\boldsymbol{r}}\times\boldsymbol{p}+\boldsymbol{r}\times\dot{\boldsymbol{p}}$，我们得到：
\begin{equation}
  \dot{\boldsymbol{M}}
  =\frac{2k}{c^{2}r^{3}}\boldsymbol{M}'\times\boldsymbol{M}.
  \tag{2}
\end{equation}

因为我们只对 $\boldsymbol{M}$ 的长期变化感兴趣，因此应当将这个表达式对粒子运动的周期进行平均.方便地进行这个平均的方法，是利用椭圆轨道运动的 $r$ 与时间关系的参数表达式，形如
\begin{equation*}
  r=a(1-e\cos\xi),
  \qquad
  t=\frac{T}{2\pi}(\xi-e\sin\xi)
\end{equation*}
（$a$ 和 $e$ 分别为椭圆的半长轴和偏心率；参见本教程第一卷§15）：
\begin{equation*}
  \overline{r^{-3}}=\frac{1}{T}\int_{0}^{T}\frac{\mathrm{d}t}{r^{3}}
  =\frac{1}{2\pi a^{3}}\int_{0}^{2\pi}
  \frac{\mathrm{d}\xi}{(1-e\cos\xi)^{2}}
  =\frac{1}{a^{3}(1-e^{2})^{3/2}}.
\end{equation*}

因此 $\boldsymbol{M}$ 的长期变化由如下公式给出：
\begin{equation}
  \frac{\mathrm{d}\boldsymbol{M}}{\mathrm{d}t}
  =\frac{2k\boldsymbol{M}'\times\boldsymbol{M}}
  {c^{2}a^{3}(1-e^{2})^{3/2}},
  \tag{3}
\end{equation}
即矢量 $\boldsymbol{M}$ 绕中心物体的旋转轴转动，其大小保持固定.

对矢量 $\boldsymbol{A}$ 的类似计算给出：
\begin{equation*}
  \dot{\boldsymbol{A}}
  =\frac{2k}{c^{2}r^{3}}\boldsymbol{M}'\times\boldsymbol{A}
  +\frac{6k}{c^{2}mr^{5}}
  (\boldsymbol{M}\cdot\boldsymbol{M}')
  (\boldsymbol{r}\times\boldsymbol{M}).
\end{equation*}

这个表达式的平均按以前一样的方式进行.从对称性的考虑可以预先判断出，平均后的矢量 $\overline{\boldsymbol{r}/r^{5}}$ 将沿着椭圆的长轴，即沿着矢量 $\boldsymbol{A}$ 的方向.这个计算导出了矢量 $\boldsymbol{A}$ 长期变化的如下表达式：
\begin{equation}
  \frac{\mathrm{d}\boldsymbol{A}}{\mathrm{d}t}
  =\boldsymbol{\Omega}\times\boldsymbol{A},
  \qquad
  \boldsymbol{\Omega}
  =\frac{2k\boldsymbol{M}'}{c^{2}a^{3}(1-e^{2})^{3/2}}
  \left\{\boldsymbol{n}'-3\boldsymbol{n}(\boldsymbol{n}\cdot\boldsymbol{n}')
  \right\}
  \tag{4}
\end{equation}
（$\boldsymbol{n}$ 和 $\boldsymbol{n}'$ 是沿 $\boldsymbol{M}$ 和 $\boldsymbol{M}'$ 的单位矢量），即矢量 $\boldsymbol{A}$ 以角速度 $\boldsymbol{\Omega}$ 旋转，而大小保持固定；最后这一点显示，轨道的偏心率不会有任何长期变化.

公式 (3) 可以写为形式
\begin{equation*}
  \frac{\mathrm{d}\boldsymbol{M}}{\mathrm{d}t}
  =\boldsymbol{\Omega}\times\boldsymbol{M},
\end{equation*}
$\boldsymbol{\Omega}$ 与 (4) 式中相同.换言之，$\boldsymbol{\Omega}$ 是椭圆“整体”旋转的角速度.这个转动既包括轨道近日点额外（与§101中考虑的相比）的移动，也包括轨道平面绕物体轴的长期转动（如果轨道面与物体的赤道面重合，就没有后面这个效应）.

为了比较起见，我们指出，对于§101中考虑的效应，相应地有
\begin{equation*}
  \Omega=\frac{6\pi km'}{c^{2}a(1-e^{2})T}\,\boldsymbol{n}.
\end{equation*}

\end{xiti}
\section{二级近似下物体系统的运动方程}

后面（§110）我们将看到，运动的物体系统要辐射引力波，因而会损失能量.这个损失只是在 $1/c$ 的第五级近似下才表现出来.在前四级近似中，系统的能量保持恒定.由此可知，不存在电磁场时，引力物体系统可以用精确到 $1/c^{4}$ 阶项的拉格朗日函数描述，而一般情形下，拉格朗日函数只精确到二阶项（§65）.这里我们来推导物体系统精确到二阶项的拉格朗日函数.这样我们就求得比牛顿近似更高一级近似的系统运动方程.

我们将忽略物体的大小和内部结构，把它们看做“类点的”粒子；换言之，在按物体的尺度 $a$ 与其相互间距 $l$ 之比的幂展开时，我们仅限于零级近似.

为了解决我们的问题，必须首先在这同级近似下，决定物体在远比其尺度大，同时又比系统辐射的引力波波长 $\lambda$ 小的距离处产生的弱引力场（$a\ll r\ll\lambda\sim lc/v$）.

精确到 $1/c^{2}$ 阶的项，远离物体的场由前节获得的表达式给出，在那里记为 $h_{ik}^{(1)}$；这里我们用的这些表达式形如 (105.6).在§105中，隐含的假定是场只由位于坐标原点的一个物体产生.但因为场 $h_{ik}^{(1)}$ 是线性化爱因斯坦方程的解，叠加原理对它有效.因此远离物体系统的场可以通过把每个物体的场简单相加求得；我们把该场写为形式
\begin{equation}
  h_{\alpha}^{\beta}=-\frac{2}{c^{2}}\varphi\delta_{\alpha}^{\beta},
\end{equation}

\begin{equation}
  h_0^0=\frac{2}{c^{2}}\varphi,\qquad h_0^{\alpha}=0,
\end{equation}
式中
\begin{equation*}
  \varphi(\boldsymbol{r})
  =-k\sum_a\frac{m_a}{|\boldsymbol{r}-\boldsymbol{r}_a|}
\end{equation*}
是点状物体系统的牛顿引力势（$\boldsymbol{r}_a$ 是质量为 $m_a$ 的物体的径矢）.度规张量为 (106.1)—(106.2) 的线元表达式为：
\begin{equation}
  \mathrm{d}s^{2}=\left(1+\frac{2}{c^{2}}\varphi\right)c^{2}\mathrm{d}t^{2}
  -\left(1-\frac{2}{c^{2}}\varphi\right)
  (\mathrm{d}x^{2}+\mathrm{d}y^{2}+\mathrm{d}z^{2}).
\end{equation}

我们注意到，含有 $\varphi$ 的一阶项不仅出现在 $g_{00}$ 中，也出现在 $g_{\alpha\beta}$ 中；在§87中已经说过，在粒子的运动方程中，$g_{\alpha\beta}$ 中的修正项给出比来自 $g_{00}$ 的项更高阶的量；因此通过同牛顿运动方程比较，我们只能确定 $g_{00}$.

我们随后将会看到，为了得到要求的运动方程，知道由 (106.1) 给出的空间分量 $h_{\alpha\beta}$ 到精度（$\sim1/c^{2}$）就够了；混合分量（在 $1/c^{2}$ 的近似中不存在）需要到 $1/c^{3}$ 阶的项，而时间分量 $h_{00}$ 需要含 $1/c^{4}$ 阶的项.为了计算它们，我们再次回到一般的引力场方程，并考虑这些方程中相应阶的项.

在不考虑物体的尺度时，我们应当把物质的能量动量张量写为形式 (33.4)，(33.5).在曲线坐标中，这个表达式重新写为
\begin{equation}
  T^{ik}=\sum_a\frac{m_ac}{\sqrt{-g}}
  \frac{\mathrm{d}x^{i}}{\mathrm{d}t}
  \frac{\mathrm{d}x^{k}}{\mathrm{d}t}
  \delta(\boldsymbol{r}-\boldsymbol{r}_a)
\end{equation}
（关于因子 $1/\sqrt{-g}$ 的出现，参见 (90.4) 中类似的过渡）；求和遍历系统中所有的物体.

分量
\begin{equation*}
  T_{00}=\sum_a\frac{m_ac^{3}}{\sqrt{-g}}\,g_{00}^{2}
  \frac{\mathrm{d}t}{\mathrm{d}s}\,\delta(\boldsymbol{r}-\boldsymbol{r}_a)
\end{equation*}
在（伽利略度规 $g_{ik}$ 的）一级近似中等于 $\sum_a m_ac^{2}\delta(\boldsymbol{r}-\boldsymbol{r}_a)$；在下一级近似中，$g_{ik}$ 由 (106.3) 代替，经简单计算后得到：
\begin{equation}
  T_{00}=\sum_a m_ac^{2}\left(1+\frac{5\varphi_a}{c^{2}}
  +\frac{v_a^{2}}{2c^{2}}\right)\delta(\boldsymbol{r}-\boldsymbol{r}_a),
\end{equation}
式中 $\boldsymbol{v}$ 是普通的三维速度（$v^{\alpha}=\mathrm{d}x^{\alpha}/\mathrm{d}t$），而 $\varphi_a$ 是场在点 $\boldsymbol{r}_a$ 的势（我们暂时不对 $\varphi_a$ 中含有的无穷大部分——粒子 $m_a$ 自场的势——加以关注；关于这一点，见下）.

至于能量动量张量的分量 $T_{\alpha\beta},T_{0\alpha}$，对它们来说，在这个近似中，只要保留 (106.4) 展开式中领头的项就够了：
\begin{equation}
\begin{aligned}
  &T_{\alpha\beta}
  =\sum_a m_a v_{a\alpha}v_{a\beta}\,\delta(\boldsymbol{r}-\boldsymbol{r}_a),\\
  &T_{0\alpha}=-\sum_a m_ac\,v_{a\alpha}\,\delta(\boldsymbol{r}-\boldsymbol{r}_a).
\end{aligned}
\end{equation}

下面来计算张量 $R_{ik}$ 的分量.利用公式 $R_{ik}=g^{lm}R_{limk}$ 来作这个计算是很方便的，式中 $R_{limk}$ 由 (92.1) 给出.这里必须记住，$h_{\alpha\beta}$ 和 $h_{00}$ 不含低于 $1/c^{2}$ 阶的项，而 $h_{0\alpha}$ 不含低于 $1/c^{3}$ 阶的项；对 $x^{0}=ct$ 微分将量小的程度提高一阶.

$R_{00}$ 中的主项为 $1/c^{2}$ 阶；除此之外我们还须保留接下去非零的 $1/c^{4}$ 阶项.简单的计算给出结果：
\begin{equation*}
\begin{aligned}
  R_{00}={}&\frac{1}{c}\frac{\partial}{\partial t}
  \left(\frac{\partial h_0^{\alpha}}{\partial x^{\alpha}}
  -\frac{1}{2c}\frac{\partial h_{\alpha}^{\alpha}}{\partial t}\right)
  +\frac{1}{2}\Delta h_{00}
  +\frac{1}{2}h^{\alpha\beta}
  \frac{\partial^{2}h_{00}}{\partial x^{\alpha}\partial x^{\beta}}\\
  &-\frac{1}{4}\left(\frac{\partial h_{00}}{\partial x^{\alpha}}\right)^{2}
  -\frac{1}{4}\frac{\partial h_{00}}{\partial x^{\beta}}
  \left(2\frac{\partial h_{\beta}^{\alpha}}{\partial x^{\alpha}}
  -\frac{\partial h_{\alpha}^{\alpha}}{\partial x^{\beta}}\right).
\end{aligned}
\end{equation*}

在这个计算中还没有对量 $h_{ik}$ 使用任何附加条件.利用这个自由，我们现在施加条件
\begin{equation}
  \frac{\partial h_0^{\alpha}}{\partial x^{\alpha}}
  -\frac{1}{2c}\frac{\partial h_{\alpha}^{\alpha}}{\partial t}=0,
\end{equation}
其结果是 $R_{00}$ 中所有含 $h_{0\alpha}$ 的项都去掉了.在余下的项中代入
\begin{equation*}
  h_{\alpha}^{\beta}=-\frac{2}{c^{2}}\varphi\delta_{\alpha}^{\beta},
  \qquad
  h_{00}=\frac{2}{c^{2}}\varphi
  +O\left(\frac{1}{c^{4}}\right),
\end{equation*}
准确到需要的精度，得到
\begin{equation}
  R_{00}=\frac{1}{2}\Delta h_{00}
  +\frac{2}{c^{4}}\varphi\Delta\varphi
  -\frac{2}{c^{4}}(\nabla\varphi)^{2}.
\end{equation}

这里我们已经转到三维表示.在计算分量 $R_{0\alpha}$ 时，保留到第一个非零阶 $1/c^{3}$ 阶的项就够了.以类似的方式，我们得到：
\begin{equation*}
  R_{0\alpha}
  =\frac{1}{2c}\frac{\partial^{2}h_{\alpha}^{\beta}}
  {\partial t\,\partial x^{\beta}}
  +\frac{1}{2}\frac{\partial^{2}h_0^{\beta}}
  {\partial x^{\alpha}\partial x^{\beta}}
  -\frac{1}{2c}\frac{\partial^{2}h_{\beta}^{\beta}}
  {\partial t\,\partial x^{\alpha}}
  +\frac{1}{2}\Delta h_{0\alpha},
\end{equation*}
然后，用条件 (106.7)：
\begin{equation}
  R_{0\alpha}=\frac{1}{2}\Delta h_{0\alpha}
  +\frac{1}{2c^{3}}\frac{\partial^{2}\varphi}
  {\partial t\,\partial x^{\alpha}}.
\end{equation}

用表达式 (106.5)—(106.9)，我们现在写出爱因斯坦方程
\begin{equation}
  R_{ik}=\frac{8\pi k}{c^{4}}
  \left(T_{ik}-\frac{1}{2}g_{ik}T\right).
\end{equation}

方程 (106.10) 的时间分量给出：
\begin{equation*}
  \Delta h_{00}
  +\frac{4}{c^{4}}\varphi\Delta\varphi
  -\frac{4}{c^{4}}(\nabla\varphi)^{2}
  =\frac{8\pi k}{c^{4}}\sum_a m_ac^{2}
  \left(1+\frac{5\varphi_a}{c^{2}}
  +\frac{3v_a^{2}}{2c^{2}}\right)
  \delta(\boldsymbol{r}-\boldsymbol{r}_a);
\end{equation*}

用恒等式
\begin{equation*}
  4(\nabla\varphi)^{2}=2\Delta(\varphi^{2})-4\varphi\Delta\varphi
\end{equation*}
和牛顿势的方程
\begin{equation}
  \Delta\varphi=4\pi k\sum_a m_a\,\delta(\boldsymbol{r}-\boldsymbol{r}_a),
\end{equation}
我们重新把这个方程写为形式
\begin{equation}
  \Delta\left(h_{00}-\frac{2}{c^{4}}\varphi^{2}\right)
  =\frac{8\pi k}{c^{2}}\sum_a m_a
  \left(1+\frac{\varphi_a'}{c^{2}}
  +\frac{3v_a^{2}}{2c^{2}}\right)\delta(\boldsymbol{r}-\boldsymbol{r}_a).
\end{equation}

在完成所有的计算后，我们已将 (106.12) 右边的 $\varphi_a$ 换为
\begin{equation*}
  \varphi_a'=-k\sum_b'\frac{m_b}{|\boldsymbol{r}_a-\boldsymbol{r}_b|},
\end{equation*}
即除物体 $m_a$ 外所有物体产生的场在 $\boldsymbol{r}_a$ 点的势.物体无穷大自势的去除（在我们把物体看做类点的方法中），相应于将它们的质量“重正化”，其结果是，它们获得了自己的真值，计及了物体本身产生的场\footnote{实际上，如果总共只有一个静止物体，方程右边就只有 $(8\pi k/c^{2})m_{a}\delta(\boldsymbol{r}-\boldsymbol{r}_{a})$，而这个方程就正确地决定了（在二级近似下）该物体产生的场.}.

利用熟悉的关系式 (36.9)：
\begin{equation*}
  \Delta\frac{1}{r}=-4\pi\delta(\boldsymbol{r}),
\end{equation*}
可以立刻得到 (106.12) 的解.于是我们有：
\begin{equation}
  h_{00}=\frac{2\varphi}{c^{2}}+\frac{2\varphi^{2}}{c^{4}}
  -\frac{2k}{c^{4}}\sum_a\frac{m_a\varphi_a'}
  {|\boldsymbol{r}-\boldsymbol{r}_a|}
  -\frac{3k}{c^{4}}\sum_a\frac{m_av_a^{2}}
  {|\boldsymbol{r}-\boldsymbol{r}_a|},
\end{equation}

方程 (106.10) 的混合分量给出：
\begin{equation}
  \Delta h_{0\alpha}
  =-\frac{16\pi k}{c^{3}}\sum_a m_av_{a\alpha}\,
  \delta(\boldsymbol{r}-\boldsymbol{r}_a)
  -\frac{1}{c^{3}}\frac{\partial^{2}\varphi}
  {\partial t\,\partial x^{\alpha}},
\end{equation}

这个线性方程的解是\footnote{在稳态情形下，方程 (106.14) 右边第二项不存在. 在离物体远距离处，它的解通过与方程 (43.4) 的解 (44.3) 类比，可以立刻写出 $h_{0\alpha}=\dfrac{2k}{c^{3}r^{2}}(\boldsymbol{M}\times\boldsymbol{n})_{\alpha}$（式中 $\boldsymbol{M}=\int\boldsymbol{r}\times\mu\boldsymbol{v}\,\mathrm{d}V=\Sigma m_{a}\boldsymbol{r}_{a}\times\boldsymbol{v}_{a}$ 是系统的角动量），与公式 (105.19) 一致.}
\begin{equation*}
  h_{0\alpha}
  =\frac{4k}{c^{3}}\sum_a\frac{m_av_{a\alpha}}
  {|\boldsymbol{r}-\boldsymbol{r}_a|}
  -\frac{1}{c^{3}}\frac{\partial^{2}f}{\partial t\,\partial x^{\alpha}},
\end{equation*}
式中 $f$ 是如下辅助方程的解
\begin{equation*}
  \Delta f=\varphi
  =-\sum_a\frac{km_a}{|\boldsymbol{r}-\boldsymbol{r}_a|}.
\end{equation*}

利用关系 $\Delta r=2/r$，我们得到：
\begin{equation*}
  f=-\frac{k}{2}\sum_a m_a|\boldsymbol{r}-\boldsymbol{r}_a|,
\end{equation*}
然后，经过简单的计算，我们得到：
\begin{equation}
  h_{0\alpha}
  =\frac{k}{2c^{3}}\sum_a
  \frac{m_a}{|\boldsymbol{r}-\boldsymbol{r}_a|}
  \left[7v_{a\alpha}
  +(\boldsymbol{v}_a\cdot\boldsymbol{n}_a)n_{a\alpha}\right],
\end{equation}
式中 $\boldsymbol{n}_a$ 是沿矢量 $\boldsymbol{r}-\boldsymbol{r}_a$ 方向的单位矢量.

表达式 (106.1)，(106.13) 和 (106.15) 足以在二阶项的精度下计算要求的拉格朗日函数.

在其他物体产生并且假设是已知的引力场中，单个物体的拉格朗日函数是
\begin{equation*}
  L_a=-m_ac\frac{\mathrm{d}s}{\mathrm{d}t}
  =-m_ac^{2}\left(1+h_{00}
  +2h_{0\alpha}\frac{v_a^{\alpha}}{c}
  -\frac{v_a^{2}}{c^{2}}
  +h_{\alpha\beta}\frac{v_a^{\alpha}v_a^{\beta}}{c^{2}}\right)^{1/2},
\end{equation*}
展开平方根并略去不重要的常数项 $-m_ac^{2}$，我们按所需精度将这个表达式重新写为
\begin{equation}
\begin{aligned}
  L_a={}&\frac{m_av_a^{2}}{2}
  +\frac{m_av_a^{4}}{8c^{2}}\\
  &-m_ac^{2}\left(\frac{h_{00}}{2}
  +h_{0\alpha}\frac{v_a^{\alpha}}{c}
  +\frac{1}{2c^{2}}h_{\alpha\beta}v_a^{\alpha}v_a^{\beta}
  -\frac{h_{00}^{2}}{8}
  +\frac{h_{00}}{4c^{2}}v_a^{2}\right).
\end{aligned}
\end{equation}

这里所有的 $h_{ik}$ 值取在点 $\boldsymbol{r}_a$；我们必须再次去掉变为无穷大的项，这相当于将作为 $L_a$ 中系数出现的质量 $m_a$ “重整化”.

进一步的计算过程如下.系统的总拉格朗日函数 $L$ 当然不等于个别物体拉格朗日函数 $L_a$ 之和，但应当这样来构造，在其他物体运动已知的情形下，使它得出作用于每个物体上力 $\boldsymbol{f}_a$ 的正确值.为此目的，我们通过微分拉格朗日函数 $L_a$ 来计算力 $\boldsymbol{f}_a$：
\begin{equation*}
  \boldsymbol{f}_a
  =\left(\frac{\partial L_a}{\partial\boldsymbol{r}}\right)
  _{\boldsymbol{r}=\boldsymbol{r}_a}
\end{equation*}
（微分是对 $h_{ik}$ 表达式中“场点”的巡行坐标 $\boldsymbol{r}$ 进行的）.然后构造总拉格朗日函数 $L$ 就容易了，由此通过求偏导数 $\partial L/\partial\boldsymbol{r}_a$ 就得到所有的力 $\boldsymbol{f}_a$.

略去简单的中间计算，我们直接给出拉格朗日函数的最后结果\footnote{相应于这个拉格朗日函数的运动方程首先是由 A. Einstein, L. Infeld, B. Hoffmann (1938) 和 A. Eddington, G. Clark (1938) 得到的.}：
\begin{equation}
\begin{aligned}
  L={}&\sum_a\frac{m_av_a^{2}}{2}
  +\sum_a\sum_b'\frac{3km_am_bv_a^{2}}{2c^{2}r_{ab}}
  +\sum_a\frac{m_av_a^{4}}{8c^{2}}
  +\sum_a\sum_b'\frac{km_am_b}{2r_{ab}}\\
  &-\sum_a\sum_b'\frac{km_am_b}{4c^{2}r_{ab}}
  \left[7(\boldsymbol{v}_a\cdot\boldsymbol{v}_b)
  +(\boldsymbol{v}_a\cdot\boldsymbol{n}_{ab})
  (\boldsymbol{v}_b\cdot\boldsymbol{n}_{ab})\right]\\
  &-\sum_a\sum_b'\sum_c'
  \frac{k^{2}m_am_bm_c}{2m_ar_{ab}r_{ac}},
\end{aligned}
\end{equation}
式中 $r_{ab}=|\boldsymbol{r}_a-\boldsymbol{r}_b|$，$\boldsymbol{n}_{ab}$ 是沿方向 $\boldsymbol{r}_a-\boldsymbol{r}_b$ 的单位矢量，求和符号上带撇的意思是应当去掉 $b=a$ 或 $c=a$ 的项.

\begin{xiti}

\noindent{\heiti 习题1}\quad 求牛顿近似中引力场的作用量.

解：用来自 (106.3) 的 $g_{ik}$，我们从一般公式 (93.3) 得到 $G=2(\nabla\varphi)^{2}/c^{4}$，所以场的作用量是
\begin{equation*}
  S_{\mathrm{g}}=-\frac{1}{8\pi k}\iint(\nabla\varphi)^{2}
  \,\mathrm{d}V\,\mathrm{d}t.
\end{equation*}

场加空间密度分布为 $\mu$ 的物质的总作用量是：
\begin{equation}
  S=\iint\left[\frac{\mu v^{2}}{2}-\mu\varphi
  -\frac{1}{8\pi k}(\nabla\varphi)^{2}\right]
  \mathrm{d}V\,\mathrm{d}t.
  \tag{1}
\end{equation}

容易验证，$S$ 对 $\varphi$ 的变分给出泊松方程 (99.2)，这是理所当然的.

能量密度由拉格朗日函数密度 $\varLambda$（(1) 式中的被积函数）用通式 (32.5) 求得，后者在目前情形下化为第二和第三项变号（因为 $\varLambda$ 中不含 $\varphi$ 对时间的导数）.将能量密度在全空间积分，以 $\mu\varphi=\varphi\Delta\varphi/(4\pi k)$ 代入第二项并作分部积分，最后得到场加物质的总能量形如
\begin{equation*}
  \int\left[\frac{\mu v^{2}}{2}
  -\frac{1}{8\pi k}(\nabla\varphi)^{2}\right]\mathrm{d}V.
\end{equation*}

因此牛顿理论中引力场的总能量是
$W=-\displaystyle\int\dfrac{(\nabla\varphi)^{2}}{8\pi k}\,\mathrm{d}V$.\footnote{为了避免任何误解，我们说，这个表达式并不同于能动赝张量（用 (106.3) 式的 $g_{ik}$ 计算）的分量 $(-g)t_{00}$；$(-g)T_{ik}$ 对 $W$ 也有贡献.}

\noindent{\heiti 习题2}\quad 求在二级近似下引力物体系统惯性中心的坐标.

解：鉴于引力相互作用的牛顿定律和电磁相互作用的库仑定律形式上完全类似，惯性中心的坐标由如下公式给出：
\begin{equation*}
\begin{aligned}
  &\boldsymbol{R}
  =\frac{1}{\mathcal{E}}\sum_a\boldsymbol{r}_a
  \left(m_ac^{2}+\frac{p_a^{2}}{2m_a}
  -\frac{km_a}{2}\sum_b'\frac{m_b}{r_{ab}}\right),\\
  &\mathcal{E}
  =\sum_a\left(m_ac^{2}+\frac{p_a^{2}}{2m_a}
  -\frac{km_a}{2}\sum_b'\frac{m_b}{r_{ab}}\right),
\end{aligned}
\end{equation*}
它与§65习题1得到的公式类似.

\noindent{\heiti 习题3}\quad 求质量相当的两个引力物体轨道近日点的长期移动 (H. Robertson, 1938).

解：两个物体所构成系统的拉格朗日函数是
\begin{equation*}
\begin{aligned}
  L={}&\frac{m_1v_1^{2}}{2}+\frac{m_2v_2^{2}}{2}
  +\frac{km_1m_2}{r}
  +\frac{1}{8c^{2}}(m_1v_1^{4}+m_2v_2^{4})\\
  &+\frac{km_1m_2}{2c^{2}r}
  \left[3(v_1^{2}+v_2^{2})-7(\boldsymbol{v}_1\cdot\boldsymbol{v}_2)
  -(\boldsymbol{v}_1\cdot\boldsymbol{n})(\boldsymbol{v}_2\cdot\boldsymbol{n})\right]
  -\frac{k^{2}m_1m_2(m_1+m_2)}{2c^{2}r^{2}}.
\end{aligned}
\end{equation*}

过渡到哈密顿函数并从中消去惯性中心的运动（参见§65习题2），我们得到：
\begin{equation}
\begin{aligned}
  \mathcal{H}={}&\frac{p^{2}}{2}\left(\frac{1}{m_1}+\frac{1}{m_2}\right)
  -\frac{km_1m_2}{r}
  -\frac{p^{4}}{8c^{2}}\left(\frac{1}{m_1^{3}}+\frac{1}{m_2^{3}}\right)\\
  &-\frac{k}{2c^{2}r}\left[3p^{2}\left(\frac{m_2}{m_1}+\frac{m_1}{m_2}\right)
  +7p^{2}+(\boldsymbol{p}\cdot\boldsymbol{n})^{2}\right]
  +\frac{k^{2}m_1m_2(m_1+m_2)}{2c^{2}r^{2}},
\end{aligned}
\tag{1}
\end{equation}
式中 $\boldsymbol{p}$ 是相对运动的动量.

我们来确定动量的径向分量 $p_r$ 作为变量 $r$、参量 $M$（角动量）和 $\mathcal{E}$（能量）的函数.这个函数决定于方程 $\mathcal{H}=\mathcal{E}$（这里在二阶项中须将 $p^{2}$ 换成其来自零级近似的表达式）：
\begin{equation*}
\begin{aligned}
  \mathcal{E}={}&\frac{1}{2}\left(\frac{1}{m_1}+\frac{1}{m_2}\right)
  \left(p_r^{2}+\frac{M^{2}}{r^{2}}\right)
  -\frac{km_1m_2}{r}\\
  &-\frac{1}{8c^{2}}\left(\frac{1}{m_1^{3}}+\frac{1}{m_2^{3}}\right)
  \left(\frac{2m_1m_2}{m_1+m_2}\right)^{2}
  \left(\mathcal{E}+\frac{km_1m_2}{r}\right)^{2}\\
  &-\frac{k}{2c^{2}r}\left[3\left(\frac{m_2}{m_1}+\frac{m_1}{m_2}\right)+7\right]
  \frac{2m_1m_2}{m_1+m_2}
  \left(\mathcal{E}+\frac{km_1m_2}{r}\right)\\
  &-\frac{k}{2c^{2}r}p_r^{2}
  +\frac{k^{2}m_1m_2(m_1+m_2)}{2c^{2}r^{2}}.
\end{aligned}
\end{equation*}

进一步的计算过程与§98中所用的类似.从上面给出的代数方程求得 $p_r$ 以后，我们对积分
\begin{equation*}
  S_r=\int p_r\,\mathrm{d}r
\end{equation*}
中的变量 $r$ 作变换，使得含 $M^{2}$ 的项化为 $M^{2}/r^{2}$.然后用小的相对论修正展开平方根下的表达式，我们得到：
\begin{equation*}
  S_r=\int\sqrt{A+\frac{B}{r}
  -\left(M^{2}-\frac{6k^{2}m_1^{2}m_2^{2}}{c^{2}}\right)\frac{1}{r^{2}}}
  \,\mathrm{d}r
\end{equation*}
（参见 (101.6)），式中 $A$ 和 $B$ 是不必具体计算出来的常系数.

结果我们得到了相对运动轨道近日点的移动：
\begin{equation*}
  \delta\varphi=\frac{6\pi k^{2}m_1^{2}m_2^{2}}{c^{2}M^{2}}
  =\frac{6\pi k(m_1+m_2)}{c^{2}a(1-e^{2})}.
\end{equation*}

与 (101.7) 比较我们看出，对于大小和形状给定的轨道，近日点的移动与物体在质量 $m_1+m_2$ 处于固定中心的场中的运动情形一致.

\noindent{\heiti 习题4}\quad 求在绕轴自转的中心物体引力场中作轨道运动的球陀螺的进动频率.

解：在一级近似下，这个效应是两个独立部分之和，一部分与中心对称场的非牛顿性质有关 (H. Weyl, 1923)，而另一部分与中心物体的转动有关 (L. Schiff, 1960).

第一部分由陀螺拉格朗日函数的附加项描述，相应于 (106.17) 中的第二项.我们将陀螺每个单元（质量为 $\mathrm{d}m$）的速度写为形式 $\boldsymbol{v}=\boldsymbol{V}+\boldsymbol{\omega}\times\boldsymbol{r}$，式中 $\boldsymbol{V}$ 是轨道运动的速度，$\boldsymbol{\omega}$ 是角速度，$\boldsymbol{r}$ 是质量元 $\mathrm{d}m$ 相对于陀螺中心的径矢（所以对陀螺体积的积分 $\int\boldsymbol{r}\,\mathrm{d}m=0$）.去掉与 $\boldsymbol{\omega}$ 无关的项，再忽略 $\boldsymbol{\omega}$ 的二次项，我们有：
\begin{equation*}
  \delta^{(1)}L=\frac{3km'}{c^{2}}\int
  \frac{\boldsymbol{V}\cdot\boldsymbol{\omega}\times\boldsymbol{r}}{R}
  \,\mathrm{d}m,
\end{equation*}
式中 $m'$ 是中心物体的质量，$R=|\boldsymbol{R}_0+\boldsymbol{r}|$ 是从场中心到质量元 $\mathrm{d}m$ 的距离，$\boldsymbol{R}_0$ 是陀螺惯性中心的径矢.在展开式 $1/R\approx1/R_0-\boldsymbol{n}\cdot\boldsymbol{r}/R_0^{2}$（$\boldsymbol{n}=\boldsymbol{R}_0/R_0$）中，第一项的积分为零，而第二项的积分用如下公式计算：
\begin{equation*}
  \int x_{\alpha}x_{\beta}\,\mathrm{d}m=\frac{1}{2}I\delta_{\alpha\beta},
\end{equation*}
式中 $I$ 是陀螺的转动惯量.结果我们得到：
\begin{equation*}
  \delta^{(1)}L=\frac{3km'}{2c^{2}R_0^{2}}
  \boldsymbol{M}\cdot(\boldsymbol{v}_0\times\boldsymbol{n}),
\end{equation*}
式中 $\boldsymbol{M}=I\boldsymbol{\omega}$ 是陀螺的角动量.

由于中心物体转动而在拉格朗日函数中出现的附加项也可以从 (106.17) 得到，但是，用§105习题中的公式 (1) 来计算它更为简单：
\begin{equation*}
  \delta^{(2)}L=\frac{2k}{c^{2}}\int
  \frac{\boldsymbol{M}'\cdot(\boldsymbol{\omega}\times\boldsymbol{r})
  \times\boldsymbol{R}}{R^{3}}\,\mathrm{d}m,
\end{equation*}
式中 $\boldsymbol{M}'$ 是中心物体的角动量.作展开，
\begin{equation*}
  \frac{\boldsymbol{R}}{R^{3}}\approx
  \frac{\boldsymbol{n}}{R_0^{2}}
  +\frac{1}{R_0^{3}}
  \left(\boldsymbol{r}-3\boldsymbol{n}(\boldsymbol{n}\cdot\boldsymbol{r})\right)
\end{equation*}
并进行积分，我们得到：
\begin{equation*}
  \delta^{(2)}L=\frac{k}{c^{2}R_0^{3}}
  \left\{\boldsymbol{M}\cdot\boldsymbol{M}'
  -3(\boldsymbol{n}\cdot\boldsymbol{M})(\boldsymbol{n}\cdot\boldsymbol{M}')\right\}.
\end{equation*}

因此，对拉格朗日函数的总修正是
\begin{equation*}
  \delta L=-\boldsymbol{M}\cdot\boldsymbol{\Omega},
  \qquad
  \boldsymbol{\Omega}
  =\frac{3km'}{2c^{2}R_0^{2}}\boldsymbol{n}\times\boldsymbol{v}_0
  +\frac{k}{c^{2}R_0^{3}}
  \left\{3\boldsymbol{n}(\boldsymbol{n}\cdot\boldsymbol{M}')
  -\boldsymbol{M}'\right\}.
\end{equation*}

与这个函数对应的运动方程是
\begin{equation*}
  \frac{\mathrm{d}\boldsymbol{M}}{\mathrm{d}t}
  =\boldsymbol{\Omega}\times\boldsymbol{M}
\end{equation*}
（参见§105中习题的方程 (2)）.这意味着，陀螺的角动量 $\boldsymbol{M}$ 以角速度 $\boldsymbol{\Omega}$ 进动，而大小保持不变.

\end{xiti}
